% !TEX program = xelatex
\documentclass[a4paper,12pt,fontset=none]{ctexart}
\usepackage{fontspec}
\setmainfont{Times New Roman}
\setCJKmainfont{Songti SC}
\setCJKsansfont{Songti SC}
\setCJKmonofont{Songti SC}
\usepackage{geometry}
\geometry{left=2.75cm,right=2.75cm,top=2.45cm,bottom=2.45cm,headsep=0.65cm,footskip=1.1cm}
\usepackage{graphicx}
\graphicspath{{../images/}{figures/}}
\usepackage{booktabs,tabularx,array,longtable,multirow}
\usepackage[table]{xcolor}
\usepackage{enumitem,fancyhdr,titlesec,indentfirst,amsmath,float,caption,hyperref,url,tikz}
% 比赛提交规范：报告内所有文字统一使用纯黑色；颜色仅保留在插图图形元素中。
\definecolor{DeepBlue}{RGB}{0,0,0}
\definecolor{AcademicBlue}{RGB}{0,0,0}
\definecolor{AccentRed}{RGB}{0,0,0}
\newcommand{\DocDate}{2026年8月31日}
\newcommand{\DocTeam}{Justifying}
\newcommand{\DocGroup}{企业赛题组}
\newcommand{\projectlink}[2]{\href{#1}{\textcolor{AcademicBlue}{\textbf{#2}}}}
\hypersetup{hidelinks,pdftitle={ESASR 项目报告},pdfauthor={Justifying}}
\setlength{\parindent}{2em}
\setlength{\parskip}{0pt}
\linespread{1.22}
% 保持赛事原模板目录层级：目录仅展示一级、二级标题。
\setcounter{tocdepth}{2}
\setlength{\headheight}{16pt}
\renewcommand{\arraystretch}{1.2}
\setlist[itemize]{leftmargin=2.2em,itemsep=0.15em,topsep=0.2em}
\setlist[enumerate]{leftmargin=2.2em,itemsep=0.22em,topsep=0.2em}
\captionsetup{font=small,labelfont=bf,labelsep=quad}
\pagestyle{fancy}
\fancyhf{}
\fancyhead[C]{\small ESASR：证据状态驱动的自适应学术检索框架}
\renewcommand{\headrulewidth}{0.6pt}
\cfoot{\small\thepage}
\fancypagestyle{plain}{\fancyhf{}\fancyhead[C]{\small ESASR：证据状态驱动的自适应学术检索框架}\renewcommand{\headrulewidth}{0.6pt}\cfoot{\small\thepage}}
\titleformat{\section}{\zihao{3}\bfseries\color{DeepBlue}}{\thesection}{0.8em}{}
\titleformat{\subsection}{\zihao{4}\bfseries\color{DeepBlue}}{\thesubsection}{0.7em}{}
\titleformat{\subsubsection}{\zihao{-4}\bfseries}{\thesubsubsection}{0.6em}{}
\titlespacing*{\section}{0pt}{1.05em}{0.6em}
\titlespacing*{\subsection}{0pt}{0.85em}{0.35em}
\titlespacing*{\subsubsection}{0pt}{0.65em}{0.25em}
% paragraph 为行内小标题；默认顶格会造成其后的正文看似漏缩进。
% 保持原有行内标题格式，仅将整条首行统一缩进两个汉字。
\titlespacing{\paragraph}{2em}{0.65em}{0.8em}
\renewcommand{\tabularxcolumn}[1]{m{#1}}
% 报告内所有自适应宽度表格列统一水平、垂直居中。
% \tabularxcolumn 已映射为 m 列，因此 X 类列同时获得垂直居中。
\newcolumntype{Y}{>{\centering\arraybackslash}X}

% 固定宽度的水平、垂直居中列
\newcolumntype{C}[1]{>{\centering\arraybackslash}m{#1}}

\setmainfont{Times New Roman}[
BoldFont = Times New Roman Bold
]

\setCJKmainfont{Songti SC}[
BoldFont = Songti SC Bold
]

\usepackage{etoolbox}
\usepackage{pdflscape}	
% 每个一级标题开始前清空浮动对象并另起一页
\pretocmd{\section}{\clearpage}{}{}
\begin{document}
\color{black}
\pagenumbering{Roman}
\thispagestyle{empty}
\begin{center}
  \vspace*{1.8cm}
  {\fontsize{25pt}{33pt}\selectfont\bfseries 第八届中国研究生人工智能创新大赛\par}
  \vspace{3.5cm}
  {\fontsize{28pt}{36pt}\selectfont\bfseries\color{DeepBlue} ESASR\par}
  \vspace{0.45cm}
  {\zihao{2}\bfseries 证据状态驱动的自适应学术检索框架\par}
  \vspace{0.3cm}
  {\zihao{4}\bfseries Evidence-State Adaptive Scholarly Retrieval\par}
  \vspace{1.0cm}
  {\zihao{2}\bfseries 项目报告\par}
  \vfill
  {\zihao{4}\bfseries 团队名称：\DocTeam\par}
  \vspace{0.25cm}
  {\zihao{4}\bfseries 参赛组别：\DocGroup\par}
  \vspace{0.25cm}
  {\zihao{4}\bfseries 提交日期：\DocDate\par}
  \vspace{1.0cm}
\end{center}

\newpage
\thispagestyle{plain}
{\small\tableofcontents}
\newpage
\thispagestyle{plain}
\begin{center}
  {\zihao{2}\bfseries 文档修订记录\par}
  \vspace{0.1em}
  \begin{tikzpicture}[x=1cm,y=-1cm,line width=0.4pt]
    % 官方模板尺寸：总宽 15.0 cm，总高 23.25 cm；表头 0.675 cm，29 个记录行各 0.7785 cm。
    \draw (0,0) rectangle (15,23.2515);
    \foreach \x in {1.435,4.425,6.675,9.165,12.160} {\draw (\x,0)--(\x,23.2515);}
    \draw (0,0.675)--(15,0.675);
    % 第 2--6 条记录各跨两行显示；外框总高度及整页版式保持不变。
    \foreach \i in {1,...,29} {
      \ifnum\i<11\relax
        \ifodd\i\draw (0,{0.675+0.7785*\i})--(15,{0.675+0.7785*\i});\fi
      \else
        \draw (0,{0.675+0.7785*\i})--(15,{0.675+0.7785*\i});
      \fi
    }

    \node[font=\bfseries\zihao{-4}] at (0.7175,0.3375) {序号};
    \node[font=\bfseries\zihao{-4}] at (2.9300,0.3375) {更改原因};
    \node[font=\bfseries\zihao{-4}] at (5.5500,0.3375) {版本};
    \node[font=\bfseries\zihao{-4}] at (7.9200,0.3375) {修订人};
    \node[font=\bfseries\zihao{-4}] at (10.6625,0.3375) {更改日期};
    \node[font=\bfseries\zihao{-4}] at (13.5800,0.3375) {备\quad 注};

    \foreach \row/\seq/\reason/\version/\reviser/\date/\note in {
      1/1/需求分析与范围界定/v0.1/队长/2026.05.22/形成需求与验收基线,
      2/2/总体架构与接口设计/v0.2/队员B/2026.06.12/数据流\\和接口规范,
      3/3/核心检索模块开发/v0.3/队员A/2026.07.10/查询编译\\与多源检索,
      4/4/基线消融与实验验证/v0.4/队长/2026.08.08/数据划分\\和评测协议,
      5/5/系统集成与材料完善/v0.5/队员B/2026.08.26/联调回归\\与材料校验,
      6/6/最终提交定版/v1.0/队长/2026.08.31/最后审阅与\\微修文档叙述
    }{
      \pgfmathsetmacro{\yc}{0.675+0.5*0.7785}
      \ifnum\row=1\relax\else\pgfmathsetmacro{\yc}{0.675+(2*\row-2)*0.7785}\fi
      \node[font=\fontsize{8.5pt}{9pt}\selectfont] at (0.7175,\yc) {\seq};
      \node[font=\fontsize{8.5pt}{9pt}\selectfont] at (2.9300,\yc) {\reason};
      \node[font=\fontsize{8.5pt}{9pt}\selectfont] at (5.5500,\yc) {\version};
      \node[font=\fontsize{8.5pt}{9pt}\selectfont] at (7.9200,\yc) {\reviser};
      \node[font=\fontsize{8.5pt}{9pt}\selectfont] at (10.6625,\yc) {\date};
      \node[font=\fontsize{8.2pt}{9pt}\selectfont,align=center] at (13.5800,\yc) {\note};
    }
  \end{tikzpicture}
\end{center}
\newpage
\pagenumbering{arabic}

\section*{摘要}
\addcontentsline{toc}{section}{摘要}

复杂科研查询往往同时包含主题、方法、数据集、时间、发表平台、开放获取和排除条件，传统关键词检索依赖反复试词，通用大模型则可能遗漏约束、缺少可核验来源并产生不稳定调用成本。本项目将其定义为\textbf{预算约束下的证据集合构建问题}：系统根据当前证据状态联合判断约束满足程度、是否继续检索以及结果集规模。

项目提出\textbf{证据状态驱动的自适应学术检索框架}（Evidence-State Adaptive Scholarly Retrieval, ESASR）。框架以约束感知查询编译、证据缺口驱动多源检索、证据融合排序、Cross Encoder 局部精排与置信度门控动态截断构成统一决策闭环；OpenAlex 与 Semantic Scholar 提供可核验候选，排序质量、绝对与相对相关性及相邻分数落差共同决定返回范围。相较固定轮次和固定 Top-K，ESASR 使检索深度与结果规模均具有可解释状态依据。

\textbf{排序模块实验}从2,536条查询中分层抽取200条，在570,206篇论文的 ScholarGym 冻结镜像上评测。等权词法召回的 Top-1 Macro F1 为0.2466；引入Cross Encoder 精排后提升至\textbf{0.2649}，相对增益约为7.42\%。在\textbf{输出策略实验}中，同一批冻结候选进一步划分为100条开发查询和100条独立测试查询：开发集仅用于确定自适应截断参数，测试集用于最终评价。自适应输出策略在测试集上的 Macro F1 为\textbf{0.2571}，高于固定 Top-1 的0.2475，平均每条查询返回1.07篇论文。\textbf{在线效率评测}从 ScholarGym 的三个查询来源中分层抽取30条有效查询，关闭规划缓存，并使用 DeepSeek-V4-Flash-0731（默认思考模式）通过真实 API 执行查询规划，同时接入 OpenAlex 与 Semantic Scholar 完成在线检索。Token 消耗直接读取模型 API 返回的 \texttt{usage} 字段，解析失败与重试产生的消耗均计入统计。30条查询中，22条进入模型规划，8条由确定性规则完成；总消耗为\textbf{24,840 Token}，按全部查询计算平均为\textbf{828 Token/条}。检索阶段共执行166次逻辑检索调用，平均为5.53次/条，未超过每条查询8次的预算上限。

工程系统基于 Vue 3、FastAPI、PostgreSQL、Redis、Neo4j 与 Docker Compose 构建，形成从复杂查询输入、多源候选组织到证据化结果输出的完整运行闭环。系统提供五档风险—覆盖策略，根据候选分数质量与相邻落差动态确定结果集规模；输出内容包括论文列表、约束覆盖摘要、检索轨迹和论文关系图，并完整保留数据来源、缺失字段、约束命中及异常状态。后端自动化测试、CLI 测试、前端类型检查与生产构建均已通过。在 MacBook Air（Apple M3，macOS 26.6.2）上开展的200条离线查询评测中，冻结候选条件下排序链路的每查询\textbf{平均处理时延为741.7 ms}，其中 Cross Encoder 精排平均耗时为285.3 ms。上述结果验证了 ESASR 的系统可运行性、过程可审计性和实验可复现性。

\textbf{关键词：}学术文献检索；证据状态决策；约束感知查询；证据缺口；多源融合；交叉编码精排；自适应结果集

\clearpage
\section{项目概况}

\subsection{背景和基础}

在选题论证、开题调研和相关工作梳理等科研任务中，研究者通常需要从海量文献中定位同时满足主题、方法、数据集、时间范围、发表平台及排除条件的论文集合。此类检索并非简单的关键词匹配，而是涉及研究意图理解、复合约束解析、子查询分解、多源迭代检索、候选质量控制、细粒度相关性排序及结果结构化组织等连续决策。现有关键词检索系统难以完整表达复杂约束，基于大语言模型的搜索方法虽然提升了语义理解和查询规划能力，但仍面临约束遗漏、检索覆盖率与精确率失衡、结果权威性与多样性难以兼顾，以及 API 调用与 Token 成本不稳定等问题。因此，复杂学术查询需要一种能够根据当前证据状态动态调整检索策略、控制资源预算并输出可核验论文集合的端到端方法。

复杂学术检索的核心矛盾并非关键词不足，而是复合研究意图难以稳定映射为可核验的论文集合。以“寻找 2021 年以后使用 Transformer 改善低纹理区域半稠密匹配、并发表于计算机视觉主会的开放获取论文”为例，主题、方法、困难场景、时间、Venue 与开放获取条件同时存在，任一约束遗漏都会造成“语义相关但任务不可用”。针对这一难点，ESASR 将自然语言意图编译为可计算约束，并以候选证据状态控制检索、排序和停止决策，使检索目标由模糊的文本相似性转化为可审计的约束满足问题。

传统检索通常将相关性、约束满足与资源开销分开处理，因而难以解释“为何继续检索”或“为何在此处截断”。ESASR 将查询形式化为 $q=(I,H,B)$：$I$ 表示主题、方法和数据集等软意图，$H$ 表示年份、排除词和强制 Venue 等硬约束，$B$ 表示检索任务、模型推理与时延预算。给定多源候选集合 $D$，系统选择规模可变的论文集合 $S_q$，在硬约束成立的条件下联合优化集合质量与调用成本：
\begin{equation}
\max_{S_q\subseteq D}\;F_1(S_q,G_q)-\lambda_1 C_{\mathrm{api}}-\lambda_2 C_{\mathrm{token}}-\lambda_3 C_{\mathrm{latency}},
\end{equation}
其中，\(G_q\) 表示评测数据集中为查询 \(q\) 标注的相关论文集合，\(S_q\) 表示系统返回的论文集合。该目标函数将候选检索、继续检索决策和结果规模选择统一为预算约束下的证据集合优化问题；相较固定检索轮次和固定 Top-K，其能够在同一框架下刻画集合质量、资源消耗与停止条件，并为后续模块设计和等预算对照实验提供统一的形式化依据。

\subsection{场景和价值}

科研选题、开题调研、相关工作（Related Work）梳理与跨学科研究通常需要同时处理主题、方法、数据集、时间范围、发表平台、开放获取及排除条件。传统关键词检索难以完整表达这些复合约束，研究者往往需要在多个平台之间反复修改检索词、人工去重并逐篇核验来源，不仅增加检索成本，也使查询过程和结果筛选依据难以复现。尤其在跨学科场景中，同义术语、元数据缺失和平台排序差异还会进一步放大漏检与误检风险。

针对上述问题，\textbf{ESASR 将复杂学术检索定义为预算约束下的证据集合构建过程}。系统首先将自然语言查询转换为结构化约束，区分必须满足的硬条件与用于判断相关性的软证据；随后通过 OpenAlex 与 Semantic Scholar 进行多源召回，并利用 DOI、来源标识和规范化标题完成跨源去重与证据合并。在候选组织阶段，系统根据当前证据缺口决定是否继续检索，通过多源证据融合与局部语义精排改善候选排序，并依据排序质量和相邻分数落差动态确定最终结果集规模。由此，问题拆解、跨源检索、候选去重、来源核验、结果排序和关系整理被组织为可追踪的统一闭环。

ESASR 的不是替代研究者作出学术判断，而是承担约束解析、跨源查找、重复记录合并、证据状态诊断和结果结构化等可程序化环节。系统完整保留查询计划、约束命中、未覆盖条件、数据来源、缺失字段、调用成本与异常状态，使研究者能够核验论文为何被检索、排序或排除。该能力适用于研究方向初筛、开题文献调研、相关工作补全、跨学科证据检索及课题组文献整理，有助于将研究者的注意力从反复试词和来源核对转移到论文阅读、证据判断与研究决策。当前实验主要验证了排序质量、输出策略和工程可运行性；系统能否进一步缩短真实用户的文献调研时间，仍需通过配对用户实验进行量化评估。

\begin{table}[H]
	\centering
	\caption{用户痛点、系统响应与验收证据}
	\small
	
	\begin{tabularx}{\textwidth}{C{2.8cm} Y Y}
		\toprule
		
		\textbf{用户痛点}
		& \textbf{系统响应}
		& \textbf{验收证据} \\
		
		\midrule
		
		复合查询约束容易遗漏
		& 将主题、方法、数据集、时间、Venue、开放获取和排除条件编译为统一约束模式
		& 查询计划、硬约束、软证据及未覆盖条件可见 \\
		
		跨平台结果重复且排序不一致
		& 基于 DOI、来源 ID 和规范化标题执行跨源去重，并融合多源排序证据
		& 原始来源、规范化标识、命中子查询及融合分量可追踪 \\
		
		固定检索轮次造成预算浪费或证据不足
		& 根据候选证据缺口和剩余预算决定是否触发定向补检
		& 缺口状态、补检原因、逻辑调用次数、HTTP 尝试、Token 与时延可记录 \\
		
		固定 Top-K 难以兼顾精确率与召回率
		& 根据候选质量、相对首位得分和相邻分数落差动态确定结果集规模
		& 独立测试集上的受控策略对照及五档策略自动化测试 \\
		
		模型生成结果来源难以核验
		& 返回学术数据源中的真实论文记录，不使用生成式论文条目替代检索结果
		& 来源、DOI、开放获取状态、缺失字段和外部接口异常显式呈现 \\
		
		检索结果难以复用和沉淀
		& 以论文列表、约束摘要、检索轨迹、详情页和关系图组织检索证据
		& 查询状态可恢复，论文可进入收藏或课题库，关系类型可区分 \\
		
		\bottomrule
	\end{tabularx}
\end{table}

\subsection{所需支持}

ESASR 已具备本地部署和基本检索能力：系统接入 OpenAlex 或 Semantic Scholar 中至少一种学术搜索 API 即可完成论文召回，在未配置大语言模型服务时，也可由确定性规则规划器完成基础查询编译。为进一步开展大规模在线评测、跨学科泛化验证和真实用户测试，项目仍\textbf{需要学术数据服务、评测数据和运行环境三方面支持}。

在学术数据服务方面，需要 OpenAlex 或 Semantic Scholar 提供稳定、高效且调用规则明确的 API，优先支持批量检索、可预期的访问额度、明确的限流与错误状态，并返回题名、摘要、DOI、作者、发表年份、Venue、开放获取状态及论文关系等可核验元数据。若能够同时稳定接入两个数据源，可进一步检验多源召回、跨源去重和证据融合的实际收益。考虑到在线索引持续更新，项目将记录 API 响应时间、数据源版本和原始响应，并构建本地冻结快照，以保证不同方法能够在一致候选条件下进行等预算比较。

在模型与运行环境方面，语义查询规划需要配置 DeepSeek 等模型服务的 API Key，本地 Cross Encoder 精排则需要具备稳定 CPU 或 GPU 推理能力的计算环境。系统已提供规则规划和基础排序降级路径，因此模型服务或精排组件不可用时仍可完成基本检索；稳定服务器和网络环境主要用于并发压力测试、在线时延评测、故障恢复测试及长时间运行验证。

在评测与用户验证方面，需要构建覆盖不同学科、不同约束复杂度和不同目标论文数量的人工标注查询集，并明确相关论文集合、硬约束满足情况和标注一致性。进一步的真实场景验证需要具有选题调研、开题报告或相关工作撰写任务的高校科研用户参与，以比较 ESASR 与常规检索流程在任务完成时间、有效论文率、约束遗漏率、来源核验成本和用户修正次数等指标上的差异。上述支持能够降低数据源波动、样本分布偏差和实验环境差异对结论的影响，并推动现有工程可行性证据扩展为可重复、可比较的外部有效性证据。

\subsubsection{当前成果边界}

ESASR 定位于基于题录、摘要与来源元数据的论文级检索和可解释证据组织，核心目标是将复杂自然语言查询转换为可执行约束，并形成来源可核验、过程可追踪、规模可自适应的论文候选集合。\textbf{当前可验收能力包括结构化查询规划、OpenAlex 与 Semantic Scholar 双源召回、跨源去重与证据融合、约束覆盖诊断、证据缺口补检、本地 Cross Encoder 精排、自适应结果集选择，以及论文详情、检索轨迹、关系图、收藏与课题库等用户工作流。系统同时提供 API、CLI 和 Docker Compose 部署入口，并统一记录数据来源、调用成本、缺失字段和外部服务异常。}

在大语言模型接入方面，ESASR 采用可配置的模型服务接口，由用户提供相应平台的 API Key。目前配置层支持 DeepSeek、Qwen、OpenAI、Kimi 以及其他兼容 OpenAI 接口规范的自定义模型服务。用户可通过 CLI 或环境变量配置模型平台、API Key、模型名称和服务地址，无需修改检索流程代码。需要强调的是，多模型接口兼容表示系统具备模型服务替换能力，并不代表所有模型均已完成同等规模的性能验证；当前正式在线 Token 与检索效率实验仅使用 DeepSeek-V4-Flash-0731，其他模型主要完成了配置与接口层面的兼容支持。

大语言模型在系统中仅承担复杂查询的语义解析和结构化规划，将自然语言输入转换为主题、方法、数据集、时间、Venue、开放获取及排除条件等统一约束，不直接生成论文条目，也不参与离线候选论文相关性指标的计算。检索结果来自 OpenAlex 和 Semantic Scholar 等真实学术数据源，本地 Cross Encoder 负责高相关候选的语义精排。因此，系统效果不归因于大语言模型生成论文答案，而归因于约束编译、证据缺口路由、多源证据融合、局部精排和自适应结果集决策共同构成的检索闭环。

受专用查询规划标注数据、训练算力和跨领域评测条件限制，当前项目不具备对基础大语言模型开展大规模预训练或专项微调的条件，也未将模型微调、参数高效训练或私有化大模型部署纳入成果范围。项目采用现有模型 API 是在研发成本、调用灵活性和结果可复现性之间作出的工程选择。模型 API Key 未配置、额度耗尽或服务不可用时，确定性规则规划器仍可完成基础查询编译，使核心检索流程不依赖单一模型平台持续可用。

在论文关系组织方面，系统仅依据学术数据源返回的关系类型构建论文网络：参考文献与施引论文表示为有向引用关系，相关论文表示为无向关联关系，不将平台推荐关系推断为真实引文、共引或因果证据。当前成果暂不覆盖独立 arXiv 数据源召回、受版权保护的全文抓取、段落级证据定位、自动综述生成及递归引文网络扩展。这些能力需要额外的数据授权、全文解析和专项评测，因此作为后续扩展方向，不纳入现阶段性能结论。

上述边界使模型使用方式、检索功能、实验指标和可运行系统保持一致，避免将模型接口兼容表述为多模型性能优势，也避免将题录与摘要检索外推为全文理解或知识生成能力。同时，统一模型服务接口、规则降级路径和结构化检索协议为后续开展模型对比、查询规划微调及全文证据抽取保留了明确的扩展入口。

\section{项目规划}

\subsection{整体目标}

固定动作链通常以预设检索轮次、固定候选规模和统一停止条件处理所有查询，难以适应不同科研问题在约束数量、表达歧义、证据稀疏程度和跨源一致性上的差异。对于主题明确、约束较少的简单查询，继续执行固定轮次容易产生不必要的模型规划与学术 API 调用；对于同时包含主题、方法、数据集、时间范围、发表平台、开放获取和排除条件的复杂查询，固定轮次又可能在关键证据尚未得到充分覆盖时提前终止。其根本原因在于，传统流程将``检索什么''、``是否继续检索''、``如何处理候选''和``最终返回多少篇论文''拆分为彼此独立的静态步骤，缺少能够描述当前证据充分程度并协调后续动作的统一状态表示。为此，ESASR 将复杂学术检索重新定义为\textbf{预算约束下的序贯证据集合构建问题}。

为统一描述检索过程中的证据变化，ESASR 定义时刻 $t$ 的证据状态为
\begin{equation}
	E_t=(C_t,G_t,S_t,B_t),
\end{equation}
其中，$C_t$ 表示当前候选集合对主题、方法、数据集、领域和开放获取等查询约束的\textbf{覆盖状态}；$G_t$ 表示尚未被有效候选覆盖的\textbf{证据缺口集合}；$S_t$ 表示候选论文经过证据融合与语义精排后形成的\textbf{有序得分分布}；$B_t$ 表示当前查询剩余的子查询数量、逻辑检索任务或外部 API 调用预算。四类状态分别回答``现有证据覆盖了什么''、``当前仍然缺少什么''、``候选之间如何区分''以及``系统还能执行多少检索动作''，从而为后续的补检、排序、动态截断与停止提供统一的决策依据。

系统首先通过约束感知查询编译器，将用户输入的自然语言科研问题转换为结构化检索条件和预算受控的子查询，并在 OpenAlex 与 Semantic Scholar 上执行首轮并行召回。候选论文返回后，系统依据 DOI、数据源标识和规范化标题完成跨源去重与记录合并，并利用论文题名、摘要和元数据更新约束覆盖状态 $C_t$，进一步得到尚未满足的证据缺口集合 $G_t$。当
\begin{equation}
	G_t\neq\emptyset
	\quad \text{且} \quad
	B_t>0.
\end{equation}
时，系统将未覆盖约束与原始研究主题重新组合为定向补检查询，并通过确定性模板触发下一轮多源召回，而不重复调用大语言模型进行查询规划。新增候选返回后，系统再次执行去重、合并和覆盖诊断，并更新证据状态 $E_{t+1}$，由此形成\textbf{``候选证据反馈—证据缺口判断—定向补检—状态更新''}的闭环。当证据缺口已经得到覆盖，或剩余任务预算不足以支持继续检索时，系统停止补检并进入候选排序阶段。与固定两轮或无条件扩展相比，该机制能够在简单查询上减少冗余调用，同时为复杂查询中尚未覆盖的关键约束提供定向补证能力。

完成候选集合构建后，多源证据融合排序模块综合词法相关性、跨来源与跨子查询的倒数排名融合、查询约束覆盖、来源一致性和权威性证据，形成可解释的基础排序。为控制深度语义模型的推理成本，本地 Cross Encoder 仅对基础排序中的\textbf{配置化 Top-$N$ 候选前缀}进行局部语义精排，并将语义相关性与基础证据分数融合，更新候选排名分布 $S_t$。在最终输出阶段，置信度门控自适应结果集策略不预设固定 Top-$K$，而是联合考察候选前缀的累计排序质量、最低绝对分数、相对首位得分比例和相邻候选分数断层，根据从``精准''到``广泛''的五档风险与覆盖策略动态确定结果截断位置。因此，``精准''档并不强制只返回一篇论文：当第二候选与首位候选近似并列且不存在明显分数断层时，系统仍可保留多篇高相关论文；``扩展''档优先覆盖主要候选质量；``广泛''档则返回候选池中达到相关性边界并满足硬约束的全部论文。该设计使最终结果规模由\textbf{查询自身的证据结构和候选得分分布}决定，而不是由预先设定的固定篇数决定。

上述机制将\textbf{补检触发、预算消耗、多源融合、局部精排、动态截断和停止条件}统一到同一证据状态框架中，使系统的每一次检索动作都具有明确的状态依据、触发条件和成本边界。与此同时，系统持续记录查询计划、触发缺口、数据来源、约束命中、缺失字段、异常状态、逻辑检索调用、Token 消耗和阶段时延，使最终输出不仅是经过排序的论文列表，而且是一个形成过程可回溯、来源信息可核验、外部故障可识别的\textbf{可审计论文证据集合}。图~\ref{fig:overall} 展示了 ESASR 如何将复杂科研问题转化为由证据状态驱动、检索深度可调、结果规模自适应且全过程可审计的论文集合。

\begin{figure}[H]
\centering
\includegraphics[width=\textwidth]{latest/Fig1_ESASR_OverallFramework.png}
\caption{ESASR 整体框架。复杂科研问题经约束编译、缺口路由、证据融合、局部精排和自适应选择形成可审计论文集合。}
\label{fig:overall}
\end{figure}

%\clearpage
%\begin{landscape}
%	\thispagestyle{plain}
%	\vspace*{\fill}
%	\begin{center}
%		\includegraphics[
%		width=0.96\linewidth,
%		keepaspectratio
%		]{latest/Fig1_ESASR_OverallFramework.png}		
%		\captionof{figure}{
%			ESASR 整体框架。复杂科研问题经约束编译、缺口路由、证据融合、
%			局部精排和自适应选择形成可审计论文集合。
%		}
%		\label{fig:overall}
%	\end{center}	
%	\vspace*{\fill}	
%\end{landscape}
%\clearpage

学术检索系统的完整性不能仅由排序指标衡量，还必须同时回答\textbf{检索结果能否复核、系统能否稳定运行，以及外部依赖失效后是否仍能保持结果可信}三个问题。针对算法实验与工程交付相互割裂的常见问题，本项目ESASR将总体目标定义为构建一套\textbf{可演示、可评测、可复现、可审计且可降级}的复杂学术检索系统。其中，“可演示”要求自然语言查询、检索规划、候选排序和结构化结果能够形成完整交互闭环；“可评测”要求排序质量、结果集选择、Token 消耗、API 调用和端到端时延具有明确口径；“可复现”要求数据划分、模型配置、评测脚本和运行环境可被重复执行；“可审计”要求论文来源、约束命中、缺失字段、检索轨迹和异常状态能够被追踪；“可降级”则要求大语言模型、学术 API 或本地精排模型不可用时，系统能够切换至规则规划、备用数据源或基础排序路径，而不是生成来源不明或状态失真的结果。

为实现上述目标，ESASR 以四项相互衔接的技术设计构成统一检索主线。\textbf{约束感知查询编译器}（Constraint-aware Query Compiler, CQC）将自然语言中的主题、方法、数据集、时间、发表平台、开放获取和排除条件转换为统一约束模式，并区分必须满足的硬约束与参与覆盖诊断和排序的软证据，从而降低复合条件遗漏和元数据不完整造成的误过滤风险。\textbf{证据缺口驱动检索路由}（Evidence-Gap-driven Retrieval Routing, EGRR）根据当前候选集合的证据覆盖、未满足约束和剩余任务预算，动态决定是否触发定向补检，使检索深度由查询实际难度决定，而不是由固定轮次决定。\textbf{多源证据融合排序}（Multi-source Evidence Fusion Ranking, MEFR）综合词法相关性、倒数排名融合、查询覆盖、来源一致性和权威性证据形成可解释基础排序；在此基础上，\textbf{局部语义精排}（Local Cross-Encoder Reranking, LCER）仅对配置化 Top-$N$ 高相关候选执行本地 Cross Encoder 推理，以较低计算成本修正词法匹配难以识别的语义误序。\textbf{置信度门控自适应结果集}（Confidence-Gated Adaptive Result Set, CARS）进一步根据候选的累计排序质量、绝对分数、相对首位比例和相邻分数断层动态确定结果截断位置，使最终返回篇数随查询证据结构变化，而不绑定固定 Top-$K$。

工程实现进一步将上述算法链路转化为可验收的系统能力。中英文交互界面用于展示查询计划、候选论文、约束覆盖、检索轨迹和论文关系；成本统计工具逐查询记录模型 Token、逻辑检索调用、底层 HTTP 尝试和阶段时延；自动化测试、命令行接口测试与容器化部署包分别验证核心逻辑、无界面运行能力和跨环境部署一致性。对于学术数据源限流或失败、模型规划输出异常、本地精排模型不可用、论文元数据缺失以及关系类型不确定等情况，系统不以补写或静默忽略掩盖异常，而是通过显式状态记录失败来源、缺失字段、降级路径和可用证据边界。由此，ESASR 将\textbf{算法设计、量化实验、运行状态与工程交付置于同一可信边界之内}：已经获得实验支持的能力以量化指标和对照结果陈述，尚未完成充分验证的功能仅作为工程机制说明，从而避免将组件数量、界面功能或技术堆叠误写为性能提升。

\subsection{技术创新点}

复杂学术检索天然涉及查询理解、候选召回、证据融合、语义排序和结果组织等多个环节。ESASR 将创新性建立在三个可检验的判据之上：其一，是否改变学术检索流程中的关键决策；其二，是否解决现有方法中具有明确边界的技术难点；其三，是否能够通过对照实验、消融分析、流程测试或运行轨迹获得验证。基于这一原则，系统不将不同组件组织为彼此独立的工具集合，而是以统一证据状态 $E_t=(C_t,G_t,S_t,B_t)$ 为主线，使查询约束、证据缺口、候选排序和剩余预算共同参与检索决策。

围绕复杂学术检索中的四个核心问题，ESASR 构建了四项相互衔接的技术创新：CQC 回答\textbf{``检索什么''}，将复合科研意图转换为可执行的硬约束与可诊断的软证据；EGRR 回答\textbf{``是否继续检索以及检索多深''}，依据证据缺口和剩余预算触发定向补检；MEFR 与 LCER 共同回答\textbf{``候选论文如何排序''}，通过多源证据融合与局部语义交叉编码精排兼顾排序质量、可解释性和推理成本；CARS 回答\textbf{``最终返回多少篇论文''}，根据候选排序质量和相邻分数结构动态确定结果截断位置。前一阶段的输出构成后一阶段的状态输入，各模块通过证据状态更新形成连续决策链，从而将查询编译、检索路由、候选排序和结果选择统一为一个逻辑闭环。

每项创新均设置明确的输入、输出、状态变化、控制变量和验证接口。已经完成受控实验的排序与结果集选择模块，以基线对比、消融实验、独立测试集和置信区间限定性能结论；查询编译与检索路由等模块，则以约束解析结果、预算上限、触发缺口、补检轨迹和自动化测试验证其机制可执行性。对于尚未完成大规模等预算在线对照的能力，报告仅将其表述为经过工程验证的系统机制，不将其外推为已经证实的质量增益。由此，ESASR 的创新性建立在\textbf{问题闭环、决策闭环与证据闭环}的统一之上。

\subsubsection{创新1：约束感知查询编译器（CQC）}

自然语言科研查询通常同时包含主题、研究方法、目标数据集、应用领域、时间范围、发表平台、开放获取和排除条件等多维信息。传统关键词检索往往将这些具有不同语义作用的条件压缩为同一关键词序列，无法区分“必须满足的任务条件”和“用于判断相关性的证据信号”，从而产生两类问题：一方面，时间、指定会议和排除项等强约束可能在查询改写过程中被遗漏，使返回论文违反用户的明确要求；另一方面，主题、方法或数据集等条件若被统一设置为硬过滤，又可能因为论文摘要缺失、术语表达不同或学术 API 元数据不完整而误删潜在相关论文。此外，若每次查询均依赖大语言模型重新理解和分解，同一请求的重复访问和并发访问还会造成额外的 Token 消耗与响应时延。为此，ESASR 提出\textbf{约束感知查询编译器}，将自然语言查询转换为结构一致、约束可执行且状态可追踪的检索计划，使后续检索、覆盖诊断、排序和审计模块共享同一语义入口。

CQC 首先建立统一约束模式，将输入查询规范化为研究问题、检索意图、结构化约束、检索子查询、歧义项和规划来源等字段。其中，结构化约束进一步划分为主题（topics）、方法（methods）、数据集（datasets）、领域（domains）、发表平台（venues）、时间下界（year\_from）、时间上界（year\_to）、开放获取要求（open access）和排除项（exclude）；同时使用 \texttt{venues\_required}、\texttt{methods\_required} 和 \texttt{primary\_topic\_required} 等标志，区分普通相关性偏好与用户明确要求必须满足的条件。统一模式的作用不仅是完成字段抽取，更重要的是为不同数据源提供稳定的中间表示：OpenAlex、Semantic Scholar、覆盖诊断模块和排序模块不需要分别解释原始自然语言，而是直接消费同一结构化计划，从而减少跨模块语义漂移，并为每项约束保留明确的来源和执行方式。

在编译流程上，CQC 采用\textbf{确定性显式约束提取与语义规划相结合的双路径机制}。系统首先通过保守的规则解析器识别查询中高置信度的年份、时间表达、会议名称、方法名称、数据集、开放获取要求和排除条件，并\textbf{将常见中文科研概念转换为 OpenAlex 与 Semantic Scholar 更易识别的英文检索表达}。对于约束结构明确、歧义较低的查询，本地规则规划器可以直接生成检索计划和子查询，不调用大语言模型；对于包含比较、排除、组合条件或复杂语义关系的查询，系统再调用配置的大语言模型执行一次 JSON 结构化规划。模型只负责将用户意图映射到预定义模式，不直接生成论文结论或虚构检索条件，并被要求返回结构化 JSON 结果和有限数量的检索子查询。由此，CQC 将大语言模型限定在其擅长的语义理解环节，而将字段校验、预算控制和约束执行保留在确定性程序中。

大语言模型生成的计划不会直接进入检索链路，而是经过规范化与确定性修复。系统将模型输出与显式约束解析结果合并，对主题、方法、数据集、领域、Venue、排除项和时间范围进行去重与类型校验；若模型遗漏了用户明确给出的年份、会议、方法或排除条件，确定性解析器会将相应字段重新写入统一计划，并通过 \texttt{constraintRepair} 记录修复是否发生、修复了哪些字段以及修复来源。年份等数值字段在进入后续模块前统一转换为合法整数，异常值被归一化为空值；重复子查询经过大小写无关的规范化去重，避免同一检索意图被多次提交给外部学术 API。该机制形成了\textbf{“语义规划负责理解复杂表达，确定性校验负责保护显式约束”}的职责边界，降低模型输出波动对检索正确性的影响。

为控制规划成本，CQC 将规范化查询、模型提供方和模型版本共同编码为缓存键，并将有效查询计划写入 Redis。相同查询再次到达时，系统直接复用缓存计划，不重复消耗模型 Token；多个相同请求并发到达且缓存尚未生成时，系统通过在途任务合并机制共享同一次规划任务，避免并发缓存未命中引发重复模型调用。运行轨迹使用 \texttt{plannerMode} 区分本地启发式规划、模型规划、缓存复用、并发合并、冻结计划复用和故障降级等状态，并分别记录模型提供方、模型版本、API 返回的 Token 使用量和规划耗时。因此，缓存命中或并发复用不会被误计为新的模型消耗，正式成本评测也可以通过无缓存规划接口绕过共享缓存，从而保证日常运行优化与实验成本计量使用不同但可核验的执行口径。

当模型平台未配置、API 请求失败、返回内容无法解析为合法 JSON，或输出结构不符合预定义模式时，CQC 不终止整个检索流程，而是切换到确定性规则规划器。降级计划与模型计划具有相同的字段结构，后续检索路由、证据融合和结果展示模块不需要因为规划来源改变而切换接口；系统同时保留 \texttt{fallbackReason}、规划模式和零模型 Token 记录，使用户能够区分“模型正常规划”和“规则降级运行”。这种同构降级机制保证的是\textbf{模型不可用时检索链路仍可执行且状态不被伪装}，而不是将规则规划的语义能力等同于大语言模型。相较于将模型异常隐藏为正常结果，显式降级更有利于界定结果可信范围和定位外部依赖故障。

上述查询编译过程如图~\ref{fig:cqc} 所示。CQC 以统一约束模式为语义接口，将确定性显式约束提取、规则规划和大语言模型语义规划组织为互补路径；随后通过字段规范化、显式约束修复、缓存复用和异常降级生成结构一致的检索计划。该流程保证了无论查询由本地规则还是大语言模型完成规划，后续检索路由、覆盖诊断与候选排序模块均能够使用相同的数据结构和约束语义

\begin{figure}[H]
	\centering
	\includegraphics[width=\textwidth]{latest/Fig2_CQC_QueryCompiler.png}
	\caption{CQC 查询编译流程。系统以统一约束模式连接确定性显式约束提取与大语言模型语义规划，并通过规范化修复、缓存复用和同构降级生成可执行的查询计划。}
	\label{fig:cqc}
\end{figure}

统一约束模式还需要解决硬约束与软证据的执行边界。若将所有条件都作为硬过滤，摘要缺失、术语变体或元数据不完整会导致潜在相关论文被过早删除；若所有条件仅参与相关性排序，返回结果又可能违反用户明确指定的时间、Venue、开放获取或排除要求。CQC 因此依据用户表达强度和字段可验证性实施分层约束：排除项对候选题名与摘要执行服务端过滤；用户明确指定的 Venue、开放获取、研究方法和核心主题分别通过对应的 \texttt{required} 状态升级为硬约束；年份边界在论文年份可判定时执行范围过滤，年份缺失则作为元数据缺失状态继续向后传播，而不以推测值补全。未被标记为必须满足的主题、方法、数据集和领域作为软证据参与候选覆盖诊断与融合排序，从而保留元数据不完整论文的后续判断机会。

对于 Venue 约束，系统不仅执行字符串匹配，还建立会议简称与正式名称之间的别名映射，例如将 CVPR、ICCV、ECCV、NeurIPS、ICLR 和 ICML 等简称与数据源返回的正式会议名称对齐。多个指定 Venue 按列表“或”语义处理，即论文命中任一允许平台即可满足该维度；当用户要求主会论文时，系统显式识别名称中的 \texttt{workshop}、\texttt{workshops} 或 \texttt{CVPRW} 等标记，防止仅因会议简称重叠而将 Workshop 论文误判为主会论文。该处理使 Venue 约束从松散关键词提升为可执行的元数据条件，同时保留关系类型和缺失字段的可追踪性。

CQC  构建了\textbf{语义理解、确定性保护、约束分层、成本复用与故障降级相结合的查询编译机制}。其输出同时服务于后续的学术 API 召回、证据缺口诊断、硬约束过滤、多源融合排序和运行轨迹展示，使自然语言意图能够沿检索链路保持一致。相较于统一关键词匹配，CQC 提高了复合约束的可执行性和检索计划的可解释性；相较于完全依赖大语言模型的查询规划，它进一步提供了显式约束修复、缓存与并发复用、Token 可计量以及无模型条件下的基本可用性。

\subsubsection{创新2：证据缺口驱动检索路由（EGRR）}

传统迭代检索通常预设固定轮次或固定查询数量，无法根据当前候选证据的充分程度调整后续动作。当首轮结果已经覆盖用户的主要约束时，继续执行固定补检会产生额外的学术 API 调用与响应时延；当首轮候选仍缺少指定方法、数据集、领域或发表平台时，固定停止又会增加关键论文漏检的风险。其根本问题在于，传统流程通常将“是否继续检索”绑定于预设轮次，而不是绑定于当前候选集合的证据状态。为此，ESASR 提出\textbf{证据缺口驱动检索路由}，以候选覆盖状态、候选规模和剩余任务预算共同决定是否进入下一轮检索，将固定动作链转换为可解释、可限制且可审计的状态决策过程。

系统首先从 CQC 生成的结构化检索计划中选取首轮子查询，并在 OpenAlex 与 Semantic Scholar 上并行召回候选论文。默认配置下，系统最多执行 4 个检索子查询，每个子查询分别提交至两个学术数据源，因此两轮检索共享最多 8 次逻辑数据源调用。首轮通常优先执行不超过 2 个核心子查询，使系统在保留后续补检空间的同时尽快形成初始候选证据。不同数据源返回的论文首先依据 DOI 构造规范键；当 DOI 缺失时，系统使用规范化标题作为跨源去重依据，仅在标题也不可用时退回数据源记录标识。对于同一规范键对应的多条记录，系统优先保留摘要信息更完整的版本，并合并命中子查询、来源平台和倒数排名信息，从而形成统一候选池。

在候选池构建后，EGRR 不以论文总数直接判断证据是否充分，而是对结构化查询中的各约束维度进行覆盖诊断。设时刻 $t$ 的候选集合为 $D_t=\{d_1,d_2,\ldots,d_n\}$，待验证约束集合为 $\mathcal{G}=\{g_1,g_2,\ldots,g_m\}$。对于候选论文 $d_i$ 和约束 $g_j$，定义维度相关的匹配函数
\begin{equation}
	m_{ij}=
	\begin{cases}
		1, & d_i \text{ 为 } g_j \text{ 提供可验证证据},\\
		0, & \text{其他情况}.
	\end{cases}
\end{equation}
进一步定义约束 $g_j$ 在候选池中的命中数量、覆盖状态和当前证据缺口集合为
\begin{equation}
	h_j=\sum_{d_i\in D_t}m_{ij},
	\qquad
	c_j=
	\begin{cases}
		1, & h_j\geq\tau_j,\\
		0, & h_j<\tau_j,
	\end{cases}
	\qquad
	G_t=\{g_j\mid c_j=0\},
\end{equation}
其中，$\tau_j$ 为不同约束维度对应的最低证据数量阈值。当前实现中，主题、方法、数据集和领域等文本证据默认至少需要 2 篇候选支持；候选文本与约束词项的重叠比例达到约束词项数量的一半时，记为一次有效文本命中。Venue、时间范围和开放获取等元数据条件采用专门的匹配逻辑，并以至少 1 条可验证记录作为覆盖依据。Venue 匹配同时执行简称与正式名称归一化，并区分主会与 Workshop；时间覆盖依据论文年份是否落入指定区间判断；开放获取覆盖则依据学术数据源返回的开放获取状态判断。由此，覆盖诊断不再依赖单一关键词是否出现，而是根据不同约束类型使用相应的证据判定规则。

首轮多源召回完成后，系统可在剩余预算允许时执行可选的引文扩展。该步骤从当前高相关候选中选择有限数量的种子论文，获取其相关论文或引用关系候选，并将新增记录重新纳入规范化去重、证据融合与覆盖诊断过程。引文扩展与关键词检索共享同一逻辑调用预算，默认最多选择 1 篇种子论文，因此不会绕过全局预算形成无界递归扩展。引文扩展并非每次查询都必须执行；当该功能未启用、缺少有效种子或剩余 API 预算不足时，系统显式将该阶段标记为跳过。

EGRR 的第二轮触发条件由\textbf{证据缺口、候选规模和剩余预算}共同确定。当引文扩展后的候选集合仍存在未覆盖约束，即 $G_t\neq\emptyset$，或当前候选数量仍低于用户请求的基本结果规模时，系统才尝试生成下一轮查询；若剩余子查询预算或逻辑数据源调用预算已经耗尽，则停止补检。对于主题、方法、数据集、领域和 Venue 等能够通过检索词表达的缺口，系统将缺口值与原始主题、核心方法或数据集组合为定向补检查询；年份和开放获取属于过滤与验证条件，不直接被机械拼接为关键词，以避免生成语义价值有限的查询。当缺口查询数量不足以利用剩余预算时，系统还可以从当前高相关候选的题名与摘要中提取重复出现、但未包含在原始查询中的术语，形成候选驱动的查询演化。缺口查询和演化查询均由确定性逻辑生成，\textbf{补检阶段不会再次调用大语言模型进行规划}。

在默认预算下，首轮查询、可选引文扩展和第二轮补检共同受最多 4 个子查询与 8 次逻辑数据源调用的上限约束。系统在生成第二轮查询前计算剩余子查询容量，在执行检索前再次计算剩余逻辑调用容量；因此，即使存在多个未覆盖条件，也不会突破全局任务预算。代码同时保留 \texttt{none}、\texttt{fixed} 和 \texttt{coverage} 三种第二轮策略：\texttt{none} 表示不执行第二轮检索，\texttt{fixed} 表示按预生成子查询执行固定补检，\texttt{coverage} 表示根据覆盖缺口和候选状态生成补检查询。这三个入口为后续开展“无补检、固定补检与缺口补检”的等预算消融实验提供了统一控制变量。

上述路由过程如图~\ref{fig:egrr} 所示。首轮 OpenAlex 与 Semantic Scholar 并行召回首先形成候选证据状态，可选引文扩展在共享预算内补充相关候选；系统随后根据约束覆盖、候选规模和剩余预算决定是否触发定向补检。补检结果重新进入去重、融合和覆盖诊断环节，形成\textbf{“多源召回—候选融合—覆盖诊断—缺口补检—状态更新”}的闭环。

\begin{figure}[H]
	\centering
	\includegraphics[width=0.6\textwidth]{latest/Fig3_EGRR_RetrievalRouting.png}
	\caption{EGRR 检索路由。首轮多源召回与可选引文扩展形成候选证据状态，约束覆盖、候选规模与剩余预算共同决定是否触发定向补检；首轮检索、引文扩展和第二轮补检共享统一任务预算。}
	\label{fig:egrr}
\end{figure}

EGRR 将\textbf{检索深度从预设轮次改写为证据状态与预算共同约束的决策变量}。对于覆盖充分且候选数量满足要求的查询，系统可以跳过第二轮检索；对于存在可表达证据缺口的查询，系统围绕缺失条件进行定向补检，而不是无差别扩展全部关键词；对于预算不足或外部数据源异常的情况，系统停止补检并保留当前候选及异常状态。每次路由均记录首轮查询数量、来源返回数量、引文扩展状态、覆盖维度、触发缺口、生成的缺口查询与演化查询、逻辑调用次数、数据源失败信息和阶段耗时，使“为什么继续检索”和“为什么停止检索”均可由运行轨迹复核。

\subsubsection{创新3：多源证据融合与局部语义精排（MEFR+LCER）}

单一学术数据源的原始排序通常由平台内部相关性模型决定，难以同时反映论文与查询的文本匹配程度、论文是否被多个子查询重复召回、不同数据源是否形成一致证据以及候选论文的引用影响力。与此同时，OpenAlex 与 Semantic Scholar 的返回分数和排序尺度并不统一，直接比较不同来源的原始分数缺乏明确语义。若进一步对全部候选执行深度语义模型推理，又会使计算成本随候选规模线性增长，放大本地部署的时延和资源压力。针对上述问题，ESASR 提出\textbf{多源证据融合排序}与\textbf{局部交叉编码精排}组成的分层排序机制：前者利用低成本、可追踪的多源证据形成基础排序，后者仅在高相关候选前缀上执行细粒度语义判断，从而在排序质量、解释能力和推理开销之间形成明确分工。

在进入 MEFR 前，系统首先执行硬约束过滤。超出用户指定时间范围的可判定记录、未命中强制 Venue 的记录、不满足开放获取要求的记录、违反排除条件的记录，以及未满足强制主题或方法要求的记录，不参与后续融合评分。该顺序保证硬约束不会被较高的引用量或来源排名抵消，也避免相关性分数掩盖结果对用户明确条件的违反。通过硬约束的候选再依据 DOI、规范化标题和来源标识进行去重，并合并其命中子查询、来源平台和来源内排名，形成统一候选表示。

对于候选论文 $d$ 和原始查询 $q$，MEFR 综合五类已归一化证据，得到基础相关性分数
\begin{equation}
	b(d,q)
	=
	0.48R(d,q)
	+
	0.25L(d,q)
	+
	0.12Q(d,q)
	+
	0.10A(d)
	+
	0.05V(d),
	\label{eq:mefr}
\end{equation}
其中，各分量均被限制在 $[0,1]$ 区间，且当前固定权重之和为 1。$R(d,q)$ 表示跨数据源与跨子查询的归一化倒数排名融合分数；$L(d,q)$ 表示查询词项在论文题名与摘要合并文本中的词项覆盖程度；$Q(d,q)$ 表示论文被不同检索子查询共同召回的比例；$A(d)$ 表示 OpenAlex 与 Semantic Scholar 等数据源对该候选形成的一致性证据；$V(d)$ 表示根据论文引用次数经过对数压缩得到的引用权威性信号。当前权重为工程实现中的固定权重，不将其表述为数据驱动学习所得的最优参数。

倒数排名融合用于消除不同数据源原始分数不可直接比较的问题。对于候选论文 $d$，系统累计其在不同数据源与不同子查询结果列表中的排名贡献：
\begin{equation}
	R_{\mathrm{raw}}(d)
	=
	\sum_{(s,u)\in\mathcal{R}(d)}
	\frac{1}{60+\operatorname{rank}_{s,u}(d)},
\end{equation}
其中，$s$ 表示数据源，$u$ 表示检索子查询，$\mathcal{R}(d)$ 表示所有召回论文 $d$ 的来源--查询组合。系统再以当前候选池中的最大倒数排名融合值完成归一化，使 $R(d,q)\in[0,1]$。因此，被多个子查询或多个数据源稳定召回、且在各列表中排名靠前的论文能够获得更高的融合贡献，而不会依赖 OpenAlex 与 Semantic Scholar 是否使用相同的原始评分尺度。

词法信号 $L(d,q)$ 由结构化查询中的研究问题、主题、方法和数据集共同形成查询词项集合，再与候选论文题名和摘要中的规范化词项计算交集比例。该信号关注查询核心术语是否能够在论文记录中获得直接文本支持，但不把词频重复出现误写为更高相关性。子查询覆盖信号 $Q(d,q)$ 则根据论文命中的不同检索子查询数量计算，并对分母进行上限控制，防止查询分解数量本身无条件放大分数。来源一致性信号 $A(d)$ 根据候选被多少个独立学术数据源共同召回计算；当前双源配置下，同时由 OpenAlex 与 Semantic Scholar 返回的论文获得更高一致性分数。引用权威性信号 $V(d)$ 使用对数压缩降低极高引用量论文对排序的支配作用，其实现形式为
\begin{equation}
	V(d)
	=
	\min\left(
	1,\,
	\frac{\log\!\left(1+\operatorname{cite}(d)\right)}
	{\log(10001)}
	\right).
\end{equation}
需要强调的是，当前实现中的 $V(d)$ 仅由引用次数计算；Venue 和开放获取状态用于约束过滤或结果说明。

MEFR 在输出基础分数的同时，为每篇论文保留规范键、来源列表、命中子查询、命中词项、基础相关性分数和元数据缺失状态。系统还根据基础分数区分“高度相关”与“部分相关”候选，并生成由命中词项、子查询数量和数据源数量组成的推荐依据。由此，基础排序不仅提供候选顺序，还保留了“论文为什么进入候选池、由哪些查询召回、得到哪些数据源支持”的解释接口，为后续错误分析、模块消融和运行审计提供可追踪中间量。

尽管 MEFR 能够融合词项匹配、来源排名和多源一致性，它仍难以充分识别同义表达、任务关系和复杂方法语义。例如，两篇论文可能共享大量术语，但实际解决不同研究问题；另一篇论文可能与查询高度相关，却因为题名和摘要使用不同表达而获得较低词法分数。为修正这类局部误序，LCER 将原始科研查询与候选论文的“题名加摘要”构造成文本对，使用本地 Cross Encoder 进行联合语义编码。模型原始输出被转换到 $[0,1]$ 区间，得到语义相关性分数 $r(d,q)$，并与 MEFR 基础分数融合为
\begin{equation}
	s(d,q)
	=
	0.65\,r(d,q)
	+
	0.35\,b(d,q),
	\label{eq:lcer}
\end{equation}
其中，$s(d,q)$ 为进入后续自适应结果集选择模块的最终排序分数。语义分数占较高权重，用于修正词法排序难以识别的语义关系；基础证据分数仍保留 0.35 的权重，以维持多源召回、查询覆盖和来源一致性等可解释证据对最终排序的约束。

LCER 不对全部候选执行深度推理，而是仅处理 MEFR 基础排序中的配置化 Top-$N$ 候选前缀。运行时实际精排数量不会超过当前候选池规模，同时至少覆盖用户请求的基本结果规模。该局部精排策略将 Cross Encoder 的计算集中于最有可能改变最终排序和结果截断位置的候选，避免在明显低相关论文上消耗计算资源。Cross Encoder 采用惰性加载机制，只有首次执行精排时才加载本地模型，并通过可配置的批大小、输入长度和计算设备适配不同部署环境。

当 Cross Encoder 未启用、依赖未安装、模型加载失败或推理过程异常时，系统不会将整个检索请求标记为失败，而是保留 MEFR 基础排序作为降级结果，同时记录精排状态、模型名称、异常类型和实际精排数量。该设计使语义精排成为可替换的质量增强层，而不是检索链路的单点依赖。MEFR 与 LCER 的分层关系如图~\ref{fig:ranking} 所示：多源召回结果首先经过硬约束过滤、规范键去重与基础证据融合，再将高相关候选前缀提交给本地 Cross Encoder，最终融合基础证据分数和语义分数形成排序结果。

\begin{figure}[H]
	\centering
	\includegraphics[width=\textwidth]{latest/Fig4_MEFR_LCER_RankingPipeline.png}
	\caption{MEFR 与 LCER 分层排序流程。候选论文在硬约束过滤后，由 MEFR 综合倒数排名融合、词项匹配、子查询覆盖、来源一致性和引用权威性形成基础分数；LCER 对配置化候选前缀执行本地 Cross Encoder 精排，并融合语义分数与基础分数。模型未启用或推理异常时，系统保留 MEFR 基础排序作为降级结果。}
	\label{fig:ranking}
\end{figure}

MEFR 与 LCER 共同构成由低成本基础证据排序到局部深度语义修正的分层处理路径。MEFR 为每篇候选保留倒数排名融合、词项匹配、子查询覆盖、来源一致性和引用权威性等分量，使基础排序能够被追踪和拆分；LCER 则以配置化 Top-$N$ 候选前缀限制深度推理范围，并在模型不可用时保留 MEFR 排序作为降级结果。两类模块分别提供可解释的全局候选压缩与细粒度的局部语义校正。

\subsubsection{创新4：置信度门控自适应结果集（CARS）}

固定 Top-$K$ 隐含了“所有查询具有相同相关论文规模”的静态假设，但复杂学术查询的候选分数结构并不一致。当多篇论文与查询近似同等相关时，固定较小的 $K$ 容易截断近似并列候选，降低证据覆盖；当候选分数在首位后快速衰减时，固定较大的 $K$ 又会引入明显低相关论文，降低结果集合的精确性。单纯根据绝对分数设置统一阈值同样难以适应不同查询的分数尺度。为此，ESASR 提出\textbf{置信度门控自适应结果集}，将最终返回篇数建模为由候选排序质量、相对首位比例、累计质量和相邻分数断层共同决定的查询相关截断问题，使结果规模随当前查询的候选证据结构动态变化，而不绑定固定 Top-$K$。

设经过硬约束过滤、MEFR 基础排序与可选 LCER 精排后的候选序列为
\begin{equation}
	s_1\geq s_2\geq\cdots\geq s_n\geq 0,
\end{equation}
其中，$s_i$ 表示第 $i$ 篇候选论文的最终排序分数。CARS 首先根据范围档位 $\lambda$ 对候选序列执行\textbf{相关性边界门控}。除首位候选始终保留外，其余候选必须同时满足最低绝对分数 $\alpha_\lambda$ 和相对首位比例 $\beta_\lambda$。由此得到连续合格候选前缀长度
\begin{equation}
	M_\lambda
	=
	\max\left\{
	m\in\{1,\ldots,n\}:
	s_i\geq\alpha_\lambda,\ 
	\frac{s_i}{\max(s_1,\epsilon)}\geq\beta_\lambda,\ 
	\forall i=2,\ldots,m
	\right\},
	\label{eq:cars-eligible}
\end{equation}
其中，$\epsilon$ 为防止首位分数为零而设置的极小正数。该门控采用连续前缀约束：一旦某个候选低于绝对分数或相对首位比例要求，其后的更低排名候选不再进入结果集，从而避免跳过低相关论文后重新选择更低排名候选所造成的排序不一致。

在合格候选前缀内，CARS 将非负排序分数归一化为相对质量权重：
\begin{equation}
	w_i
	=
	\frac{s_i}{\sum_{j=1}^{M_\lambda}s_j},
	\qquad
	C_k
	=
	\sum_{i=1}^{k}w_i,
	\quad
	1\leq k\leq M_\lambda,
	\label{eq:cars-mass}
\end{equation}
其中，$C_k$ 表示前 $k$ 篇候选所覆盖的累计排序质量。当合格候选的分数总和为零时，系统保留首位候选并将其质量权重设为 1，以保证空分数或异常分数条件下仍有确定性输出。需要强调的是，$w_i$ 表示\textbf{当前候选池内部的归一化排序质量}，不是经过概率校准的论文相关概率。

给定档位 $\lambda$ 的目标质量覆盖 $\tau_\lambda$，系统首先确定达到目标质量的最小基础截断位置：
\begin{equation}
	k_\lambda^{(0)}
	=
	\min\left\{
	k\leq M_\lambda:
	C_k\geq\tau_\lambda
	\right\}.
	\label{eq:cars-initial-cutoff}
\end{equation}
若在该位置直接截断，可能把与边界候选近似并列的下一篇论文排除。CARS 因而继续检查相邻候选之间的分数落差，并从 $k_\lambda^{(0)}$ 开始沿连续候选前缀向后扩展，直至出现大于档位阈值 $\delta_\lambda$ 的显著分数断层，或到达相关性边界 $M_\lambda$。最终截断位置定义为
\begin{equation}
	K_\lambda
	=
	\min\left(
	\left\{
	k\in\{k_\lambda^{(0)},\ldots,M_\lambda-1\}:
	s_k-s_{k+1}>\delta_\lambda
	\right\}
	\cup
	\{M_\lambda\}
	\right).
	\label{eq:cars-final-cutoff}
\end{equation}
该式表明，CARS 并不是寻找“同时达到质量目标且存在分数断层”的单一步骤截断点，而是先达到目标质量，再保护边界附近的近似并列候选，最终在显著分数断层或相关性边界处停止。因而即使处于“精准”档，只要第二篇及后续候选与首位候选分数接近，系统仍可返回多篇论文。

CARS 的完整工程流程如图~\ref{fig:cars} 所示。系统首先接收按最终分数降序排列的候选序列，通过绝对分数与相对首位比例构造连续合格前缀，再依据累计排序质量确定基础截断位置，并沿近似并列候选继续扩展，直至检测到显著分数断层或触及相关性边界。最终输出始终是原始排序的连续前缀，因此不会出现跳过高排名候选而保留低排名候选的结果结构。

\begin{figure}[H]
	\centering
	\includegraphics[width=0.8\textwidth]{latest/Fig5_CARS_AdaptiveSelection.png}
	\caption{CARS 自适应结果集选择流程。系统首先依据最低绝对分数和相对首位比例确定连续合格候选前缀，再根据累计归一化排序质量确定基础截断位置，并沿近似并列候选继续扩展，直至出现显著分数断层或到达相关性边界；五档策略均不绑定固定返回篇数。}
	\label{fig:cars}
\end{figure}

CARS 将“最终返回多少篇论文”从固定超参数转化为由查询内部候选分布决定的状态决策。其计算仅使用已经生成的排序分数，不增加学术 API 调用、大语言模型调用或 Cross Encoder 推理，因此能够以较低额外开销调整结果集合规模。

\section{实施方案}

\subsection{技术可行性分析}


ESASR 的技术可行性取决于三个基本条件：复杂查询能否被转化为确定且可执行的检索任务，多源论文数据与语义模型能否在受控预算内完成候选构建，以及系统能否在外部服务不稳定时保持结果真实和核心功能可用。现有实现与测试结果表明，上述条件均具备明确的工程基础。OpenAlex 与 Semantic Scholar 提供论文题名、摘要、作者、发表时间、来源标识和引用信息等基础数据，大语言模型仅承担一次结构化查询规划，本地 Cross Encoder 仅处理基础排序后的有限候选前缀；检索任务数、模型 Token、HTTP 请求和精排规模均设置显式上限。因此，系统不依赖无限轮次搜索或全量深度推理，能够在有限计算资源和可控调用成本下完成端到端检索。

为将上述算法流程落实为可运行系统，ESASR 采用分层架构隔离用户交互、检索编排、外部数据服务和状态存储。如图~\ref{fig:dataflow} 所示，前端负责复杂查询输入及证据呈现，FastAPI 负责查询规划、检索路由、候选排序和结果组织，PostgreSQL、Redis 与 Neo4j 分别承担业务数据、检索缓存和论文关系存储。结构化接口使算法模块能够独立替换，也使外部数据源、模型服务和存储组件的异常能够被分别定位。

系统运行所需资源与现有实验条件相匹配。查询规划采用远端模型 API，不要求本地部署或微调大语言模型；语义精排使用轻量级本地 Cross Encoder，并仅作用于配置化候选前缀。系统已在 MacBook Air（Apple M3，macOS 26.6.2）上完成 200 条离线查询测试，平均端到端延迟为 741.7\,ms，其中本地语义精排平均耗时为 285.3\,ms。该结果表明，核心检索与排序流程能够在普通个人计算设备上运行，不依赖专用 GPU 集群，具备现场演示和轻量化部署条件。

外部依赖失效不会被静默处理。模型服务未配置、密钥无效或规划结果异常时，系统使用确定性规则规划器生成同构查询计划；本地 Cross Encoder 不可用时，系统保留多源证据融合的基础排序；单一学术数据源超时或限流时，系统继续使用其他来源的有效结果，并在检索轨迹中记录异常来源和失败原因。后端自动化测试、CLI 测试、前端类型检查与生产构建均已通过，Docker Compose 提供数据库、缓存和服务组件的一体化启动方式。由此，ESASR 已具备可实现、可运行、可部署和可降级的技术基础。

\begin{figure}[H]
	\centering
	\includegraphics[width=\textwidth]{latest/Fig6_System_EndToEndDataFlow.png}
	\caption{系统端到端数据流与可信边界。交互层、服务编排层、学术数据服务和状态存储层通过结构化接口解耦，论文来源、字段缺失、约束命中和异常状态随检索过程显式传播。}
	\label{fig:dataflow}
\end{figure}

\subsection{技术细节}

\subsubsection{数据来源、缺失值与关系语义}

学术检索系统的可信性不仅取决于排序质量，还取决于每条论文记录能否追溯至真实数据源。外部 API 请求失败可能污染质量评测。因此，ESASR 将\textbf{来源标识、调用状态和字段完整性}作为候选证据的一部分，使数据真实性约束贯穿召回、融合、排序、展示与评测全流程。

对于 OpenAlex 与 Semantic Scholar 真实请求发生超时、限流或响应异常时，系统将失败来源、异常信息和调用状态写入检索轨迹及健康状态；其他可用数据源返回的论文仍可进入后续处理，从而避免单一来源故障扩散为全链路失败。对于外部论文记录中缺失的年份、摘要、数字对象唯一标识符（Digital Object Identifier, DOI）等元数据，系统保留原始空值，不根据标题、相邻记录或模型输出进行推测性补全，并通过 \texttt{metadataMissing} 字段向后端流程与前端界面显式暴露缺失项。排序模块仅使用实际可获得的证据计算对应分量，界面则将字段缺失与来源异常分别呈现，避免把“数据源未提供该字段”误解释为“论文不存在”或“约束未命中”。

论文关系图还面临另一类语义风险，即把平台生成的相关论文推荐误解释为可核验的引文关系。为避免不同关系类型在图结构中被错误合并，ESASR 对引文关系和相关性关系采用独立编码。对于参考文献和施引文献（citing papers），系统统一使用有向关系 \texttt{CITES}，并规定边 $u\rightarrow v$ 表示论文 $u$ 引用了论文 $v$：当前论文的参考文献形成“当前论文指向被引论文”的边，施引文献则形成“施引论文指向当前论文”的边。该方向定义与学术引用事实保持一致，可用于回答“该论文引用了谁”和“哪些论文引用了该论文”两类问题。

OpenAlex 返回的 related works 仅表示平台计算得到的主题或内容相关性，不能证明两篇论文之间存在直接引用。系统因此将其编码为无向关系 \texttt{RELATED}，并保留关系来源，不将其参与真实引文方向的推断。前端使用带箭头的实线表示 \texttt{CITES}，使用无箭头的虚线表示 \texttt{RELATED}，同时在关系详情中展示关系类型及数据来源。通过从数据获取、后端存储到前端呈现的一致语义约束，系统既保留了相关论文探索的辅助价值，又防止用户将统计相关性误判为可核验的引用事实，从而提高关系图的可解释性与学术严谨性。

\subsubsection{部署、安全与可降级运行}

多依赖检索系统的工程风险不仅来自单个组件失效，还包括密钥进入前端产物、健康状态与真实服务能力不一致，以及局部故障被扩大为整次检索失败。ESASR 因此将部署目标定义为\textbf{密钥不越过服务端边界、运行状态能够被外部编排系统判断、可替代能力优先降级而非整体中断}，并通过容器隔离、分级健康检查和请求级降级机制实现上述目标。

如图~\ref{fig:deploy} 所示，Docker Compose 将 Web、API、PostgreSQL、Redis 与 Neo4j 划分为独立服务。Web 容器仅负责静态页面分发和请求反向代理，不保存模型服务或学术检索服务的访问密钥；大语言模型与学术搜索 API 的密钥仅注入 API 容器，真实配置以只读方式挂载，不写入前端构建产物。PostgreSQL、Redis 和 Neo4j 分别保存业务数据、缓存及任务状态和论文关系，外部用户不直接访问这些状态服务。由此，用户流量、服务端密钥和内部业务状态在部署层形成明确边界，降低前端源码泄漏、配置误发布和内部服务暴露的风险。

\begin{figure}[H]
	\centering
	\includegraphics[width=\textwidth]{latest/Fig9_Deployment_Topology.png}
	\caption{ESASR 容器化部署拓扑。用户请求经 Web 容器转发至 API 服务，模型与学术数据源访问、检索编排和内部状态读写均在服务端完成；访问密钥、业务数据与关系状态不进入前端构建产物。}
	\label{fig:deploy}
\end{figure}

系统将\textbf{部署就绪状态}与\textbf{检索能力状态}分别管理。健康接口检查 API 服务自身以及 PostgreSQL、Redis 和 Neo4j 等内部依赖；当不可替代的核心依赖异常时，接口返回 HTTP 503，并将总体状态标记为 \texttt{overall=degraded}，使 Docker Compose 或上层部署平台能够依据真实状态判断实例是否就绪。模型规划服务、单一学术数据源和本地精排模型具有替代路径，其异常原则上不直接终止整个检索请求，而是触发能力降级，并将切换原因、失败来源和调用尝试写入运行轨迹。该分级机制避免把“部分高级能力不可用”与“系统核心状态不可用”混为同一故障。

具体降级策略如表~\ref{tab:failure-degradation} 所示。模型规划服务未配置、密钥无效或输出无法解析时，系统切换至确定性规则规划器，继续生成与语义规划器结构一致的查询计划；单一学术 API 触发限流或请求失败时，系统执行受控退避，并在预算允许时保留另一来源的有效结果；本地 Cross Encoder 未安装或推理异常时，系统跳过局部语义精排，继续使用多源证据融合的基础排序。上述策略均以\textbf{不伪造结果、不隐藏异常、不突破预算}为约束，使系统在能力下降时仍能返回边界明确且可核验的结果。

\begin{table}[H]
	\centering
	\caption{关键故障、降级行为与可信性保证}
	\label{tab:failure-degradation}
	\begin{tabularx}{\textwidth}{C{3.0cm} Y Y}
		\toprule
		\textbf{故障情形} & \textbf{系统行为} & \textbf{可信性保证} \\
		\midrule
		模型规划服务未配置或输出异常
		& 切换至确定性规则规划器，生成同构查询计划
		& 查询核心流程仍可执行，切换原因进入运行轨迹 \\
		
		单一学术 API 限流或请求失败
		& 执行全局限速和受控退避，并在预算允许时使用另一来源
		& 失败来源和 HTTP 尝试可见 \\
		
		本地 Cross Encoder 不可用
		& 跳过局部语义精排，保留多源证据融合的基础排序
		& 高级排序能力降级，但不扩大为整次检索失败 \\
		
		PostgreSQL 在开发环境中不可达
		& 通过显式开发配置切换至本地 SQLite
		& 降级仅限开发与功能验证，不将 SQLite 状态伪装为生产环境状态 \\
		
		不可替代的核心依赖异常
		& 健康接口返回 HTTP 503 和 \texttt{overall=degraded}
		& 部署编排依据真实依赖状态判断实例是否就绪 \\
		\bottomrule
	\end{tabularx}
\end{table}

该部署方案将安全控制、服务编排、运行观测和降级策略纳入同一可信边界。其工程价值不在于宣称系统在任意故障下均可完整运行，而在于明确区分\textbf{正常运行、能力降级和核心不可用}三种状态：可替代模块发生异常时保留基本检索能力，不可替代依赖失效时拒绝伪装为健康实例，所有状态变化均通过健康接口或检索轨迹向外暴露。由此，系统能够避免“界面可访问即服务可用”的错误判断，并为现场演示、自动化部署和故障复核提供一致的运行语义。

\subsubsection{主题适配、多语言与交互一致性}

复杂学术检索界面同时包含查询约束、论文元数据、检索轨迹、成本指标和关系图等高密度信息。如果界面文本、组件样式和业务状态直接耦合，不同语言之间容易出现术语不一致、文本溢出和状态提示缺失，主题切换也可能破坏图表辨识度。ESASR 因此将\textbf{语言资源、视觉主题和检索业务状态}分别管理，使界面外观与显示语言的变化不改变查询语义、排序结果和审计信息。

在多语言支持方面，系统提供中文和英文两套界面资源，并以稳定的语义键统一管理导航、查询槽位、约束提示、阶段状态、异常信息和关系类型等文本。用户切换界面语言时，仅替换交互标签和系统说明，不改写原始查询，也不翻译外部数据源返回的论文题名、摘要、作者和发表平台信息，从而避免二次翻译改变学术含义。数字对象唯一标识符（Digital Object Identifier, DOI）、Venue、OpenAlex、Semantic Scholar、\texttt{CITES} 和 \texttt{RELATED} 等具有稳定语义的专业标识保持统一形式；缺失的语言资源回退至默认文本，而不是显示空白内容。由此，界面国际化与学术数据真实性相互隔离，保证同一查询在不同界面语言下对应相同的检索计划和候选集合。

界面语言与查询语言属于两个独立维度。中文界面可以提交英文或中英混合查询，英文界面也可以处理中文研究问题；查询编译模块依据输入内容识别主题、方法、数据集、时间和发表平台等约束，而不依赖当前界面语言。该设计避免将“切换显示语言”错误等同于“翻译用户查询”，并使多语言能力覆盖真实科研场景中常见的中文问题描述、英文专业术语和中英混合约束表达。

在主题适配方面，系统通过统一的设计变量管理背景、正文、边框、强调状态和交互反馈，支持浅色与深色主题之间的一致切换。主题变化仅影响视觉呈现，不修改检索参数、候选顺序和用户工作流状态；论文列表、检索轨迹、成本统计和关系图共享同一套主题状态，避免普通组件已经切换而图表仍保留原主题的问题。ECharts 与 G6 的节点、边、标签和提示框根据当前主题同步更新，同时通过箭头、虚实线和文字标签表达引文方向、相关关系与异常状态，不将关键语义仅编码为颜色，以保证不同显示环境下的信息可辨识性。

为减少长时间检索过程中的重复操作，主题和语言偏好由前端统一维护，并与查询内容、筛选条件、排序方式和滚动位置分别保存。用户切换语言或主题后，当前查询计划、已返回论文和检索轨迹保持不变；从论文详情返回结果页时，系统恢复原有筛选、排序和浏览位置，避免因界面设置变化重新调用模型或学术 API。该机制使主题、多语言和页面导航成为可恢复的交互状态，而不是触发检索链路重新执行的业务事件。

上述设计的工程价值在于建立\textbf{显示层变化不影响检索事实}的一致性边界。主题切换解决不同光照环境和长时间阅读条件下的可读性问题，多语言资源降低跨语言科研用户的操作成本，业务状态隔离则保证界面设置不会改变实验结果或增加额外调用。由此，ESASR 在保持学术数据原貌和检索过程可复核的前提下，提高了系统的可访问性、交互连续性与现场展示稳定性。

\subsection{实验设计与结果}

\subsubsection{实验问题与分层验证协议}

ESASR 同时改变查询解析、检索深度、候选排序和结果规模。若只报告完整链路的单一指标，将无法判断性能变化来自查询计划改善、补检范围扩大、候选次序调整，还是返回篇数增加。为避免不同模块之间的效果混合归因，本项目按照“\textbf{模块决策是否改变、控制变量是否明确、结论是否具有对应证据}”组织实验，并建立四个相互隔离的验证层级：

\begin{enumerate}
	\item \textbf{查询编译验收}：使用包含多类约束的固定复杂查询，检验约束感知查询编译器（CQC）能否生成可执行计划，并保持硬约束、软证据与缺失状态可追踪；
	\item \textbf{检索路由实验}：在冻结语料和统一候选条件下，比较固定补检与证据缺口驱动检索路由（EGRR），检验其能否在质量非劣条件下降低调用成本；
	\item \textbf{候选排序实验}：冻结查询、语料和 Top-100 候选池，分别检验多源证据融合排序（MEFR）中的跨路由证据，以及局部语义交叉编码精排（Local Cross-Encoder Reranking, LCER）的独立贡献；
	\item \textbf{输出策略实验}：冻结候选及其排序，仅改变结果集截断规则，比较置信度门控自适应结果集（CARS）与固定 Top-K；
	\item \textbf{在线效率实验}：接入真实模型服务、OpenAlex 与 Semantic Scholar，测量规划 Token、逻辑检索调用、外部服务异常和端到端墙钟时间。
\end{enumerate}

上述分层分别回答“查询能否正确执行”“何时继续检索”“候选顺序是否改善”“应当返回多少篇”和“真实运行需要多少资源”。其中，带 Gold Set 的冻结候选实验用于支持性能结论；配对消融用于支持模块归因；真实链路复测和自动化测试仅用于证明工程可执行性。未建立专项对照的能力不被表述为已经取得性能优势。图~\ref{fig:protocol} 展示了各实验层级的数据关系与控制边界。

\begin{figure}[H]
	\centering
	\includegraphics[width=\textwidth]{latest/Fig7_Experiment_EvidenceProtocol.png}
	\caption{ESASR五层实验与证据协议。CQC固定复杂查询验收查询编译与约束传播；EGRR在冻结语料上比较固定补检与收益门控；MEFR与LCER在统一Top-100候选池内分别检验多源证据和局部语义精排；CARS在开发集选择规则并于独立测试集比较固定Top-K；在线效率实验接入真实模型与双学术数据源测量Token、调用、时延和外部故障。五条路径分别验证查询编译、检索深度、候选排序、结果规模与运行成本，避免不同效应混合归因。}
	\label{fig:protocol}
\end{figure}

\subsubsection{数据集、评测单位与指标定义}

在线学术索引持续更新，论文元数据、默认排序和接口返回规模均可能随时间变化。若直接比较不同时间的在线响应，数据漂移可能被误判为算法改进。因此，EGRR、MEFR、LCER 与 CARS 的质量实验均使用 ScholarGym 冻结本地镜像。该镜像包含 570,206 篇论文和 2,536 条查询，查询来源覆盖 PaSa AutoScholar、PaSa RealScholar 与 LitSearch。

实验首先使用历史200条分层查询进行特征分析、参数校准和路由阈值选择，随后从剩余查询中抽取两批互不重叠、且均未参与参数选择的200条查询，分别作为首次验证集和独立复验集。CARS 实验另将一组200条冻结查询按来源和稳定哈希拆分为100条开发查询与100条独立测试查询，测试集 Gold Set 不参与阈值选择。

每条查询$q$对应一个人工标注的相关论文集合，即Gold Set，记为$G_q$。论文身份依次通过arXiv ID、DOI、数据源ID和规范化标题进行匹配；标题匹配仅在稳定标识符缺失时使用，以降低同名论文和元数据差异造成的误配风险。统计分析以\textbf{查询}为独立单位，而不是将候选论文数量视为样本量。

对于查询$q$的返回集合$R_q$，逐查询Precision、Recall与F1定义为

\begin{equation}
	P_q=\frac{|R_q\cap G_q|}{|R_q|},
	\qquad
	R_q^{\mathrm{rec}}=\frac{|R_q\cap G_q|}{|G_q|},
\end{equation}

\begin{equation}
	F1_q=
	\frac{2P_qR_q^{\mathrm{rec}}}
	{P_q+R_q^{\mathrm{rec}}}.
\end{equation}

主指标Macro F1为全部测试查询逐查询F1的算术平均：

\begin{equation}
	\operatorname{MacroF1}
	=
	\frac{1}{|Q|}
	\sum_{q\in Q}F1_q.
\end{equation}

排序实验同时报告Recall@20、MAP@20和nDCG@20。其中，Recall@20衡量前20篇候选对Gold Set的覆盖程度；MAP@20同时考虑相关论文是否出现及其排序位置；nDCG@20对位于前列的相关论文赋予更高权重。所有置信区间均以查询为重采样单位，通过5,000次非参数bootstrap估计；方法间增益采用配对bootstrap，即在每次重采样中同时计算相同查询上的两种方法。单方法区间用于描述指标不确定性，配对差值区间用于判断方法增益能否排除零差异，两者不混用。

\subsubsection{CQC 复杂约束编译验收}

复杂科研查询通常同时包含主题、方法、数据集、年份、Venue、开放获取和排除项。若全部条件被直接处理为关键词，年份、发表平台和排除词等硬约束容易在召回阶段失效。CQC 将自然语言查询映射到统一约束模式，并将必须满足的条件划分为硬约束，将用于覆盖诊断、补检与排序的条件划分为软证据。

实验采用同时包含主题、方法、困难场景、数据集偏好、时间范围、Venue、开放获取和排除项的固定查询：

\begin{quote}
	\small
	寻找2020年以来使用Transformer提升弱纹理或大视角变化场景下图像匹配性能的论文，要求发表于CVPR、ICCV或ECCV且可开放获取，优先包含ScanNet或MegaDepth实验，排除纯光流方法和综述文章。
\end{quote}

\begin{table}[H]
	\centering
	\caption{CQC固定复杂查询的链路验收结果}
	\label{tab:cqc-case}
	\small
	\begin{tabular}{
			>{\centering\arraybackslash}m{2.0cm}
			>{\centering\arraybackslash}m{1.7cm}
			>{\centering\arraybackslash}m{1.7cm}
			>{\centering\arraybackslash}m{1.7cm}
			>{\centering\arraybackslash}m{1.7cm}
			>{\centering\arraybackslash}m{1.7cm}
			>{\centering\arraybackslash}m{1.8cm}}
		\toprule
		\textbf{规划方式} &
		\textbf{解析时间/ms} &
		\textbf{候选子查询} &
		\textbf{执行子查询} &
		\textbf{原始候选} &
		\textbf{最终结果} &
		\textbf{硬约束违规} \\
		\midrule
		确定性CQC & 14 & 5 & 4 & 120 & 4 & 0 \\
		\bottomrule
	\end{tabular}
\end{table}

如表 \ref{tab:cqc-case} 确定性CQC在14\,ms内生成5个英文候选子查询，EGRR根据任务预算执行其中4个。OpenAlex与Semantic Scholar共返回120条原始记录，经规范键去重和硬约束过滤后保留4篇论文。整个链路执行8次数据源调用，端到端耗时为9.088\,s。最终返回的LoFTR、ASpanFormer、CasMTR与CoMatch均满足年份、指定会议、开放获取、Transformer方法和图像匹配主题等硬约束。系统计算得到的证据覆盖率为60\%。该实验表明，CQC能够将包含多类约束的自然语言查询转化为可执行计划，并保持约束命中与缺失证据可追踪。

\subsubsection{EGRR 检索路由质量--成本实验}

固定执行第二轮检索无法区分证据已经充分的简单查询与关键条件仍未覆盖的复杂查询，容易在前者上浪费调用预算。EGRR根据约束覆盖、证据缺口、候选规模和剩余预算估计补检的边际收益，仅在门控条件满足时执行第二轮检索。

实验在相同冻结语料、查询计划、候选规模和基础排序下比较两种策略：固定补检对所有查询执行预设第二轮；收益门控EGRR则根据当前证据状态决定是否继续检索。首次验证集用于观察路由行为；在独立复验前，将nDCG@20非劣界值固定为0.02，以检验EGRR能否在质量损失不超过可接受范围的条件下降低调用数。

\begin{table}[H]
	\centering
	\caption{EGRR的排序质量与逻辑检索调用成本}
	\label{tab:egrr-replication}
	\small
	\begin{tabular}{
			>{\centering\arraybackslash}m{2.3cm}
			>{\centering\arraybackslash}m{3.0cm}
			>{\centering\arraybackslash}m{1.8cm}
			>{\centering\arraybackslash}m{1.8cm}
			>{\centering\arraybackslash}m{2.1cm}
			>{\centering\arraybackslash}m{2.1cm}}
		\toprule
		\textbf{数据划分} &
		\textbf{检索策略} &
		\textbf{Recall@20} &
		\textbf{nDCG@20} &
		\textbf{平均逻辑调用} &
		\textbf{触发第二轮} \\
		\midrule
		验证集1（200） &
		固定补检 &
		0.3507 & 0.2394 & 3.00 & 200/200 \\
		
		验证集1（200） &
		\textbf{收益门控EGRR} &
		0.3395 & 0.2340 & \textbf{1.06} & 6/200 \\
		\midrule
		独立复验集（200） &
		固定补检 &
		0.4171 & 0.3260 & 3.00 & 200/200 \\
		
		独立复验集（200） &
		\textbf{收益门控EGRR} &
		\textbf{0.4261} & \textbf{0.3290} &
		\textbf{1.03} & 3/200 \\
		\bottomrule
	\end{tabular}
\end{table}

首次验证中，EGRR相对固定补检的nDCG@20差值为$-0.0054$，配对95\%置信区间为$[-0.0113,-0.00004]$；独立复验中的差值为$+0.0030$，区间为$[-0.0024,0.0084]$。独立复验区间的下界高于预先固定的非劣界值$-0.02$，支持EGRR在该界值下质量非劣；首次验证的结果方向与该结论一致。

两批共400条查询中，EGRR仅对9条查询触发第二轮检索，平均逻辑调用数为1.045次；固定补检平均需要3.00次，调用量减少约65.2\%。因此，\textbf{在nDCG@20满足预设非劣条件的前提下，收益门控路由减少了不必要的逻辑检索调用。} 该模块解决了是质量约束下的预算分配问题。

\subsubsection{MEFR 多源证据分量消融}

多源候选可能由不同数据源和子查询重复召回。倒数排名融合（Reciprocal Rank Fusion, RRF）能够记录论文在不同检索路径中的排名一致性，但其独立排序收益需要通过跨数据划分验证。实验在相同查询、相同Top-100候选池和相同硬约束条件下，比较校准基础排序与加入RRF证据后的结果。

\begin{table}[H]
	\centering
	\caption{MEFR中RRF分量的跨数据划分消融}
	\label{tab:mefr-rrf}
	\small
	\begin{tabular}{
			>{\centering\arraybackslash}m{2.3cm}
			>{\centering\arraybackslash}m{3.2cm}
			>{\centering\arraybackslash}m{1.7cm}
			>{\centering\arraybackslash}m{1.8cm}
			>{\centering\arraybackslash}m{1.7cm}
			>{\centering\arraybackslash}m{1.8cm}}
		\toprule
		\textbf{数据划分} &
		\textbf{排序配置} &
		\textbf{F1@1} &
		\textbf{Recall@20} &
		\textbf{MAP@20} &
		\textbf{nDCG@20} \\
		\midrule
		验证集1（200） &
		校准基础排序 &
		0.1424 & 0.3395 & 0.1885 & 0.2334 \\
		
		验证集1（200） &
		基础排序+RRF &
		0.1424 & 0.3507 & 0.1924 & 0.2394 \\
		\bottomrule
	\end{tabular}
\end{table}

如表 \ref{tab:mefr-rrf} 实验结果表明，加入校准RRF后nDCG@20提高0.0060，配对95\%置信区间为$[0.0006,0.0117]$。

\subsubsection{LCER 局部语义精排复验}

词法排序难以识别同义表达和复杂方法关系，但对全部候选执行Cross Encoder会显著增加推理成本。LCER使用本地
\texttt{cross-encoder/ms-marco-MiniLM-L6-v2}
对校准基础排序的前20篇候选进行查询--论文联合编码，其余候选保持原有顺序。两批实验共享570,206篇论文的冻结镜像和相同Top-100候选池，因此性能变化仅来自候选内部次序调整，而非候选集合变化。

\begin{table}[H]
	\centering
	\caption{LCER在两批未见查询上的独立复验结果}
	\label{tab:lcer-replication}
	\small
	\begin{tabular}{
			>{\centering\arraybackslash}m{2.3cm}
			>{\centering\arraybackslash}m{3.2cm}
			>{\centering\arraybackslash}m{1.9cm}
			>{\centering\arraybackslash}m{1.8cm}
			>{\centering\arraybackslash}m{1.7cm}
			>{\centering\arraybackslash}m{1.8cm}}
		\toprule
		\textbf{数据划分} &
		\textbf{排序方法} &
		\textbf{Top-1 Macro F1} &
		\textbf{Recall@20} &
		\textbf{MAP@20} &
		\textbf{nDCG@20} \\
		\midrule
		验证集1（200） &
		校准基础排序 &
		0.1424 & 0.3395 & 0.1885 & 0.2334 \\
		
		验证集1（200） &
		\textbf{基础排序+LCER} &
		\textbf{0.1899} & 0.3395 &
		\textbf{0.2219} & \textbf{0.2585} \\
		\midrule
		独立复验集（200） &
		校准基础排序 &
		0.2257 & 0.4274 & 0.2787 & 0.3292 \\
		
		独立复验集（200） &
		\textbf{基础排序+LCER} &
		\textbf{0.2466} & 0.4274 &
		\textbf{0.2988} & \textbf{0.3454} \\
		\midrule
		合并分析（400） &
		校准基础排序 &
		0.1841 & 0.3834 & 0.2336 & 0.2813 \\
		
		合并分析（400） &
		\textbf{基础排序+LCER} &
		\textbf{0.2183} & 0.3834 &
		\textbf{0.2604} & \textbf{0.3019} \\
		\bottomrule
	\end{tabular}
\end{table}

在首次验证集中，LCER将Top-1 Macro F1、MAP@20和nDCG@20分别提高0.0475、0.0334和0.0250，对应配对95\%置信区间分别为$[0.0179,0.0804]$、$[0.0149,0.0535]$和$[0.0102,0.0413]$。三项区间均未包含0。

独立复验中，三项指标分别提高0.0209、0.0201和0.0162，变化方向与首次验证一致，但该批查询上的单独差值区间包含0。合并两批共400条未见查询后，LCER将Top-1 Macro F1提高0.0342，区间为$[0.0095,0.0604]$；MAP@20提高0.0267，区间为$[0.0116,0.0421]$；nDCG@20提高0.0206，区间为$[0.0082,0.0330]$。三项合并区间均位于0以上，支持\textbf{LCER能够改善固定候选池内的首位命中与前列排序质量}。LCER在首次验证集和独立复验集上的平均精排耗时分别为417.8\,ms和425.2\,ms，合并平均约为421.5\,ms。

\subsubsection{CARS 自适应输出策略实验}

固定Top-K隐含所有查询具有相同答案规模的假设。较小的$K$可能遗漏近似并列论文，较大的$K$则可能引入低相关候选。CARS根据绝对排序分数、相对首位分数、累计质量和相邻分数落差确定连续候选前缀，使返回篇数随查询证据结构变化。

为隔离“返回多少”的影响，实验冻结同一批候选及其排序，并将200条查询按来源和稳定哈希拆分为100条开发查询与100条独立测试查询。两侧均包含78条AutoScholar、2条RealScholar和20条LitSearch查询。CARS仅在开发集搜索8,064组规则；测试集Gold Set不参与阈值或最大返回数选择。固定Top-1、Top-2、Top-3、Top-5与CARS共用相同的测试查询和候选排序。

\begin{table}[H]
	\centering
	\caption{CARS自适应输出策略独立测试结果（$n=100$）}
	\label{tab:cars}
	\small
	\begin{tabular}{
			>{\centering\arraybackslash}m{2.4cm}
			>{\centering\arraybackslash}m{1.5cm}
			>{\centering\arraybackslash}m{1.5cm}
			>{\centering\arraybackslash}m{1.8cm}
			>{\centering\arraybackslash}m{2.7cm}
			>{\centering\arraybackslash}m{2.1cm}}
		\toprule
		\textbf{输出策略} &
		\textbf{Precision} &
		\textbf{Recall} &
		\textbf{Macro F1} &
		\textbf{F1 95\% CI} &
		\textbf{平均返回数} \\
		\midrule
		固定Top-1 &
		0.3300 & 0.2241 & 0.2475 &
		[0.1715, 0.3255] & 1.00 \\
		
		固定Top-2 &
		0.1900 & 0.2537 & 0.1986 &
		[0.1424, 0.2537] & 2.00 \\
		
		固定Top-3 &
		0.1400 & 0.2697 & 0.1665 &
		[0.1219, 0.2105] & 3.00 \\
		
		固定Top-5 &
		0.0960 & 0.2926 & 0.1304 &
		[0.0991, 0.1627] & 5.00 \\
		
		\textbf{CARS} &
		\textbf{0.3350} &
		\textbf{0.2374} &
		\textbf{0.2571} &
		\textbf{[0.1812, 0.3350]} &
		\textbf{1.07} \\
		\bottomrule
	\end{tabular}
\end{table}

固定增大$K$提高了Recall，却持续降低Precision与Macro F1。CARS仅在7\%的测试查询上返回第二篇论文，Macro F1由固定Top-1的0.2475提高至0.2571，绝对增益为0.0097；平均返回规模仅由1.00篇增加至1.07篇。CARS相对固定Top-1的配对差值区间为$[0.0000,0.0267]$。因此，当前结果支持\textbf{轻量门控呈现正向的质量--规模平衡趋势}。

\subsubsection{真实联网规划与检索效率}

冻结语料上的零Token仅说明排序实验隔离了远端服务成本，不能代表完整系统的真实运行开销。在线效率实验从PaSa AutoScholar、LitSearch和PaSa RealScholar中分别无放回抽取10条有效查询，共30条，随机种子为20260828。系统采用EGRR配置接入OpenAlex与Semantic Scholar，单条查询最多执行8次逻辑检索调用。

语义规划使用DeepSeek-V4-Flash-0731，API模型名为
\texttt{deepseek-v4-flash}，启用默认思考模式并强制绕过规划缓存。Token仅根据API响应中的\texttt{usage}字段统计，不使用本地字符数估算；LCER为本地模型，不产生规划Token，其推理成本与大模型规划成本分别计量。

\begin{table}[H]
	\centering
	\caption{真实联网规划与多源检索效率（$n=30$）}
	\label{tab:online-token}
	\small
	\begin{tabular}{
			>{\centering\arraybackslash}m{3.1cm}
			>{\centering\arraybackslash}m{1.8cm}
			>{\centering\arraybackslash}m{1.8cm}
			>{\centering\arraybackslash}m{2.5cm}
			>{\centering\arraybackslash}m{3.2cm}}
		\toprule
		\textbf{指标} &
		\textbf{总量} &
		\textbf{均值} &
		\textbf{中位数/P95} &
		\textbf{计量口径} \\
		\midrule
		总Token &
		24,840 &
		828.0 &
		886/2,453 &
		API \texttt{usage} \\
		
		输入Token &
		5,811 &
		193.7 &
		240.5/474 &
		30条完整查询 \\
		
		输出Token &
		19,029 &
		634.3 &
		645/1,979 &
		包含思考Token \\
		
		思考Token &
		15,759 &
		525.3 &
		482.5/1,838 &
		输出Token的子集 \\
		
		逻辑检索调用 &
		166 &
		5.53 &
		4/8 &
		单条上限8次 \\
		
		在线链路耗时/ms &
		802,290 &
		26,743 &
		15,911.5/90,065 &
		真实API墙钟时间 \\
		\bottomrule
	\end{tabular}
\end{table}

30条查询中，22条进入模型规划，8条由确定性规则完成；32次模型请求包含2次解析失败后的重试，失败尝试产生的Token全部计入成本。22条模型规划查询的平均消耗为1,129.1 Token，中位数为977，P95为2,453。30条检索流程均成功完成，平均返回19.67篇候选；其中3条查询记录到Semantic Scholar HTTP异常，对应10\%的外部故障查询率，但另一学术来源和系统降级路径保证了流程完成。

在线实验表明，规则规划能够避免部分模型调用，单条查询平均消耗828 Token和5.53次逻辑检索调用。平均墙钟时间为26.743\,s，而P95达到90.065\,s，说明真实链路的主要不确定性来自远端模型和学术API响应，而不是本地候选排序。该结果验证了系统具备可计量的在线运行闭环，同时也表明后续优化应优先面向API并发、限流处理、超时控制和缓存复用。

\subsection{作品完整性与展示闭环}

\subsubsection{系统交付与复现入口}

ESASR 通过\projectlink{https://github.com/esasr/esasr}{项目代码仓库}交付 npm CLI、Docker Compose 编排文件、配置模板、启动脚本和实验复现入口。CLI 引导用户配置模型服务与学术搜索 API Key；未配置大语言模型或模型规划异常时，系统自动切换至确定性规则规划器，保留基本检索能力并记录降级原因。

项目将查询编译、检索路由、候选排序、自适应输出和在线运行分别映射到可执行脚本与结果文件，使评委能够从同一代码仓库复现关键流程。代码、配置、命令和实验产物共同构成作品交付边界，避免将界面演示等同于完整系统。

\subsubsection{工程验证与测试覆盖}

页面能够访问并不意味着检索系统已经具备工程可用性。ESASR 因此从后端逻辑、CLI 发布、前端构建、核心模块实验和真实联网运行五个层面建立验证入口，如表~\ref{tab:verification}所示。

\begin{table}[H]
	\centering
	\caption{工程验证与实验复现入口}
	\label{tab:verification}
	\small
	\begin{tabularx}{\textwidth}{
			C{2.5cm}
			C{2.5cm}
			C{5.0cm}
			C{3.4cm}}
		\toprule
		\textbf{验证对象} &
		\textbf{验证结果} &
		\textbf{覆盖内容} &
		\textbf{复现入口} \\
		\midrule
		
		后端自动化测试 &
		63/63 通过 &
		来源、缺失元数据、约束过滤、补检、精排、动态范围、缓存、限流和成本指标 &
		\texttt{tests/} \\
		
		CLI 与发布打包 &
		9/9 通过；300.4\,kB &
		初始化、端口选择、密钥配置、服务重启、发布文件白名单和打包完整性 &
		\texttt{npm test}、\texttt{pack:check} \\
		
		前端生产构建 &
		类型检查与构建通过 &
		Vue 3、TypeScript 类型约束、组件依赖和生产静态产物 &
		\texttt{npm run build} \\
		
		CQC 查询编译验收 &
		固定复杂查询完成 &
		八类约束解析、硬软约束划分、候选子查询生成、硬约束过滤和缺失状态传播 &
		CQC 验收脚本 \\
		
		EGRR 路由复验 &
		200验证+200复验 &
		固定补检与收益门控对照、nDCG@20 非劣检验、调用数量和触发比例 &
		核心模块评测脚本 \\
		
		MEFR 与 LCER 复验 &
		200验证+200复验 &
		统一 Top-100 候选、RRF 分量消融、F1@1、MAP@20、nDCG@20和精排时延 &
		排序复验脚本 \\
		
		CARS 输出评测 &
		100开发+100测试 &
		固定 Top-K、门控规则选择、测试集隔离、平均返回数和配对 bootstrap &
		输出策略评测脚本 \\
		
		真实联网效率评测 &
		30/30 查询完成 &
		模型规划 Token、双源检索调用、墙钟时间、重试和外部服务异常 &
		在线效率评测脚本 \\
		\bottomrule
	\end{tabularx}
\end{table}

上述验证形成了从算法实现、接口测试、前端构建到真实联网运行的连续证据。自动化测试验证关键状态转换能够稳定执行，冻结实验用于评估模块质量与资源成本，真实联网评测则暴露远端服务时延和异常。三类证据相互补充但不相互替代：构建通过不等价于算法有效，固定案例成功也不等价于跨查询性能优势。

\subsubsection{可审计交互与结构化输出}

传统学术搜索界面通常只呈现最终论文列表，用户难以判断查询约束是否被正确解析、第二轮检索为何触发、候选排序依据来自何处，以及外部数据源是否发生异常。ESASR 将查询计划、首轮召回、覆盖诊断、缺口补检、多源融合、局部精排和结果集选择组织为连续阶段状态，并在结果页中同步展示约束命中、未覆盖条件、数据来源、检索轨迹和成本指标。

论文卡片呈现题名、作者、Venue、年份、摘要、来源、推荐理由与开放获取入口；详情页进一步汇总稳定标识符、来源信息和分型关系，并提供收藏与课题库入口。系统输出真实论文记录而不是生成式论文清单，字段缺失时保留空值，数据源失败时保留异常状态，避免使用默认内容制造表面完整性。

该交互将算法内部状态转化为用户可理解、可核验的外部证据，使系统形成“自然语言输入、约束编译、证据检索、状态解释、论文核验与结果沉淀”的完整工作流。界面承担过程解释与证据组织，检索质量仍由独立实验验证，而不由演示流畅度或视觉效果替代。

\subsubsection{演示查询设计}

功能演示若只采用包含少量关键词的简单查询，容易将普通文本搜索包装为复杂学术检索能力。为同时检验约束编译、预算路由、双源召回、硬约束过滤和缺失证据传播，作品展示采用以下查询：

\begin{quote}
	\small
	寻找2020年以来使用Transformer提升弱纹理或大视角变化场景下图像匹配性能的论文，要求发表于CVPR、ICCV或ECCV且可开放获取，优先包含ScanNet或MegaDepth实验，排除纯光流方法和综述文章。
\end{quote}

该查询包含主题、方法、困难场景、数据集偏好、时间范围、Venue、开放获取和排除项八类条件。其中，时间范围、指定Venue、开放获取要求和排除项形成可执行的硬约束；困难场景与数据集偏好作为软证据参与覆盖诊断与补检。该划分能够同时检验系统的约束执行能力与不完整元数据处理能力。

\subsubsection{演示链路与检索结果解释}

在2026年8月26日的本地在线复测中，确定性CQC在14\,ms内生成5个英文候选子查询，EGRR在任务预算内执行其中4个；OpenAlex与Semantic Scholar共返回120条原始记录，经规范键去重和硬约束过滤后保留4篇论文。整个链路执行8次数据源调用，不产生远端模型规划Token，端到端耗时为9.088\,s，且未发生数据源请求失败。该结果与表~\ref{tab:cqc-case}中的查询编译验收保持一致，说明复杂约束可以在不依赖模型规划的条件下进入真实检索流程。

最终结果包括CVPR 2021的LoFTR、ECCV 2022的ASpanFormer、ICCV 2023的CasMTR与ICCV 2025的CoMatch。4篇论文均满足2020年以来、指定会议、开放获取、Transformer方法和图像匹配主题等硬约束，硬约束违规数为0。系统计算得到的证据覆盖率为60\%，未能由题录或摘要稳定支持的数据集偏好与困难场景仍被标记为缺口，而没有被强制填充为已满足状态。

该演示的验收重点不是返回论文数量，而是系统能否在统一链路中同时完成约束执行、来源核验和不确定性暴露。LCER不可用时，系统保留MEFR校准基础排序并记录降级原因；单一学术数据源异常时，系统保留另一来源的真实结果和失败轨迹。因此，固定查询演示证明的是系统的\textbf{可运行性、可审计性和可降级性}，不能替代前述跨查询消融与独立复验。

\subsubsection{系统界面与展示路径}

自然语言检索首页如图~\ref{fig:home}所示。输入区支持用户表达主题、方法、时间、数据集、Venue、开放获取和排除项，并通过示例提示降低复杂查询的表达成本。

\begin{figure}[H]
	\centering
	\includegraphics[width=\textwidth]{system-home.png}
	\caption{自然语言检索首页。输入区支持主题、方法、时间、数据集、Venue、开放获取和排除条件等复杂查询表达，并提供范围策略与示例提示。}
	\label{fig:home}
\end{figure}

检索结果页如图~\ref{fig:results}所示。页面在同一视图中呈现结构化查询计划、约束覆盖状态、检索阶段轨迹、成本信息和带来源证据的论文候选，使用户能够核验查询条件是否进入实际检索流程。

\begin{figure}[H]
	\centering
	\includegraphics[width=\textwidth]{system-search-results.png}
	\caption{检索结果页。结构化查询计划、约束覆盖状态、检索轨迹、成本信息和带来源证据的论文候选在同一页面呈现。}
	\label{fig:results}
\end{figure}

论文详情页如图~\ref{fig:detail}所示。页面汇总可核验元数据、摘要、开放获取入口、来源标识和分型关系信息，并支持收藏与课题库沉淀。引用关系与平台相关性关系采用不同语义和视觉编码，避免将相关论文推荐误解释为真实引文关系。

\begin{figure}[H]
	\centering
	\includegraphics[width=\textwidth]{system-paper-detail.png}
	\caption{论文详情页。页面汇总可核验元数据、摘要、开放获取入口、来源标识和分型关系信息，并支持收藏和课题库沉淀。}
	\label{fig:detail}
\end{figure}

\subsubsection{完整性结论与能力边界}

ESASR 已形成从CLI配置、容器化启动、复杂查询输入、双源召回、状态驱动补检、候选排序、自适应输出到论文收藏和关系组织的完整交付链路。后端测试、CLI测试、前端生产构建、冻结模块实验和真实联网效率评测分别提供代码正确性、发布完整性、工程可构建性、算法有效性和在线资源成本证据。

作品完整性并不意味着所有模块均已达到相同的实验强度。CQC当前主要由固定复杂查询和流程测试支持；EGRR与LCER已经完成两批未见查询复验；CARS现有证据仅覆盖一至两篇的受控门控特例；五档风险--覆盖策略仍需更大规模专项评测。明确区分已验证性能、工程能力与后续扩展，使系统展示与实验结论共享同一可信边界。


\subsection{应用价值与科研服务前景}

\paragraph{科研应用定位。}

大规模学术索引擅长覆盖，语言模型擅长理解自然语言，但前者难以表达复合约束，后者又缺少稳定来源与成本边界。ESASR 不重复建设底层索引或生成论文清单，而是以 CQC、EGRR、MEFR、LCER 和 CARS 将两类能力组织为\textbf{可控检索编排层}，并将来源、覆盖、预算与失败状态贯穿科研工作流。该定位使系统能够复用开放学术基础设施，同时提供通用平台缺失的约束可执行性、过程可复现性和故障可降级性，适用于选题调研、Related Work、跨学科检索与课题组证据管理。

\paragraph{科研服务与技术验收。}

\begin{table}[H]
\centering
\caption{科研服务场景与技术验收指标}
\begin{tabularx}{\textwidth}{C{3.0cm} Y Y}
\toprule
\textbf{科研场景} & \textbf{系统支持} & \textbf{主要验收指标} \\
\midrule
复杂选题与开题调研 & 复合约束编译、多源召回和证据覆盖诊断 & 有效论文率、硬约束违反率和任务完成时间 \\
Related Work 与跨学科检索 & 可解释排序、局部语义精排和来源追踪 & Top-K 检索质量、来源完整性和重复结果比例 \\
课题组证据管理 & 共享课题库、检索轨迹和论文关系组织 & 结果可复现率、检索记录完整性和用户可用性 \\
科研信息服务 & 检索编排 API、统一身份与结构化结果接口 & 并发能力、p95 时延、故障恢复时间和数据合规性 \\
\bottomrule
\end{tabularx}
\end{table}

现有离线指标能够表明排序趋势，却不能直接回答系统是否减少研究者的检索负担。系统已提供账户、课题库、检索轨迹、来源核验和统一指标日志，可在开题调研与 Related Work 任务中记录完成时间、有效论文率、重复检索率、约束遗漏率和用户修正次数，并与常规检索流程进行配对比较。该用户级验证把算法质量转化为科研效率与结果可用性证据，是从技术可行性迈向科研服务价值所需的关键验收环节，而非泛化的应用前景描述。

\subsection{可靠性保障与扩展验证}

真实科研服务必须同时面对外部学术API限流、模型输出异常、本地精排资源不足、元数据缺失、状态服务故障和公网访问风险。任何单点异常都可能导致检索中断，或更隐蔽地生成来源不明、字段伪完整和状态不可复核的结果。ESASR 因此不以“零故障”为可靠性目标，而是将可靠性定义为：\textbf{故障能够被识别、影响范围能够被限制、降级行为能够被解释、恢复条件能够被验证。}

\subsubsection{学术数据源波动与多源降级}

OpenAlex与Semantic Scholar具有不同的限流策略、元数据覆盖和服务可用性，若将任一来源视为稳定依赖，单源异常便可能扩大为整次检索失败。ESASR在API访问层实施全局速率控制、请求超时、指数退避和缓存复用，并通过双源并行召回降低单一来源不可用造成的影响。重试次数、底层HTTP尝试、异常状态和失败来源均写入运行轨迹，使用户能够区分“未检索到相关论文”与“数据源未成功返回”。

当单一来源发生限流或临时异常时，系统保留另一来源返回的真实记录，并显式标注来源缺失，不使用本地演示论文补充计分结果。在30条真实联网查询中，3条记录到Semantic Scholar HTTP异常，但30条检索流程均完成，说明多源互补和故障隔离能够避免单一数据源异常直接扩大为完整链路失败。该结果只能证明当前样本中的可降级执行能力，不能外推为任意故障强度下均可保持相同检索质量。

\subsubsection{模型规划异常与确定性接管}

大语言模型能够解析复杂语义约束，但其输出可能出现字段缺失、格式错误、响应超时或服务未配置等问题。若下游模块直接依赖未经校验的模型文本，错误将传播至子查询生成、约束过滤和检索路由。CQC因此要求模型输出符合统一JSON约束模式，并在进入检索链路前执行字段类型、枚举值和范围校验。

模型服务未配置、请求失败或输出无法通过模式校验时，系统切换至确定性规则规划器，生成与语义规划器结构一致的查询计划。切换原因、实际规划路径、重试次数和Token消耗均保留在运行轨迹中。规则规划不保证与模型规划具有相同的语义解析能力，但能够维持主题、年份、Venue、开放获取、包含词和排除词等基本条件的可执行性，从而将“模型不可用”限制为规划能力降级，而不是整次检索中断。

\subsubsection{排序模块降级与结果一致性}

局部语义交叉编码精排需要额外模型文件、内存和推理时间，不能被视为始终可用的基础依赖。LCER未安装、模型加载失败或推理异常时，系统不返回部分精排或新旧分数混合的候选序列，而是完整保留MEFR校准基础排序，并记录精排未执行的原因。

该降级方式能够保证结果仍由词项覆盖、原始召回位置、约束命中和来源一致性等可解释证据产生。根据排序复验，LCER能够进一步改善固定候选池内的前列排序质量，因此关闭LCER可能降低排序效果；但基础排序仍可执行，来源、硬约束和候选身份不会因精排缺失而改变。由此，系统将“质量增强模块不可用”与“检索结果不可用”区分处理，避免局部资源不足扩大为全链路故障。

\subsubsection{元数据真实性与不确定性传播}

学术数据源普遍存在摘要、年份、DOI、Venue和开放获取状态缺失。若使用默认值填补未知字段，系统可能错误执行硬约束或生成无法核验的推荐理由。ESASR对来源字段采用“真实值、缺失值、异常状态”三类语义：缺失年份、摘要或DOI保持为空，并通过\texttt{metadataMissing}状态向前端传播；跨源记录依据DOI、数据源ID和规范化标题去重，合并时优先保留信息更完整的真实字段及全部来源标识。

离线展示数据与真实检索结果采用严格隔离。兼容展示接口中的离线卡片统一标记为\texttt{Offline Demo}或\texttt{offline\_fallback}，不得进入正式计分池和在线质量结论。系统也不将平台返回的相关论文自动解释为引用关系：参考文献和被引论文编码为有向\texttt{CITES}，平台相关性记录编码为无向\texttt{RELATED}。该设计以显式不确定性替代伪完整数据，使筛选、排序、关系图和用户解释共享同一事实来源。

\subsubsection{部署边界与访问安全}

公网部署的主要风险包括密钥进入前端产物、跨域配置过宽、身份令牌失效和内部状态服务直接暴露。Docker Compose将Web、API、PostgreSQL、Redis与Neo4j隔离部署：Web容器仅承载静态页面和反向代理，模型密钥与学术API Key只注入API容器，不进入浏览器代码或前端构建产物；PostgreSQL、Redis与Neo4j仅通过内部容器网络向后端开放。

访问控制采用HTTPS保护传输链路，通过CORS来源白名单限制浏览器跨域访问，并使用JWT完成用户身份校验。配置文件不写入代码仓库，真实密钥通过环境变量或只读配置注入，并支持按部署周期轮换。上述机制降低了密钥泄漏和未授权访问风险，但不能替代专业安全审计；正式公网部署仍需结合反向代理访问日志、依赖漏洞扫描和最小权限配置持续验证。

\subsubsection{健康检查与运行可观测性}

页面能够打开并不意味着检索服务已经就绪。ESASR的健康接口联合检查API、PostgreSQL、Redis与Neo4j状态；任一核心依赖异常时返回HTTP 503和\texttt{overall=degraded}，使容器编排和运维人员能够依据真实依赖状态判断服务是否可用，而不是以静态页面可访问作为就绪标准。

完整检索接口按查询记录规划方式、模型请求次数、Token消耗、逻辑检索调用、底层HTTP尝试、缓存命中、精排耗时、端到端时延、失败来源和降级原因。这些指标将算法行为、外部服务成本和工程故障分账记录，使异常能够定位到查询编译、远端检索、候选排序或状态存储等具体阶段，并为后续并发压力测试和故障注入提供统一观测口径。

\begin{table}[H]
	\centering
	\caption{关键风险、系统控制与可核验状态}
	\label{tab:reliability}
	\small
	\begin{tabularx}{\textwidth}{
			C{3.0cm}
			C{4.7cm}
			C{4.4cm}
			C{2.2cm}}
		\toprule
		\textbf{风险场景} &
		\textbf{控制与降级行为} &
		\textbf{可观测证据} &
		\textbf{验证状态} \\
		\midrule
		
		单一学术API限流或异常 &
		全局限速、超时、指数退避、缓存复用，并保留另一数据源结果 &
		失败来源、HTTP尝试、重试次数、返回来源和端到端时延 &
		流程与在线验证 \\
		
		模型服务未配置或输出异常 &
		统一模式校验失败后切换确定性规则规划器 &
		规划路径、切换原因、模型请求、Token和解析错误 &
		自动化测试 \\
		
		LCER模型不可用 &
		取消本次语义精排，完整保留MEFR校准基础排序 &
		精排状态、降级原因、基础分数和最终排序 &
		自动化测试 \\
		
		论文元数据缺失 &
		保持空值并传播\texttt{metadataMissing}，不使用默认值伪填充 &
		缺失字段、来源标识和跨源合并记录 &
		自动化测试 \\
		
		离线展示数据混入 &
		统一标记\texttt{Offline Demo}/\texttt{offline\_fallback}，禁止进入计分池 &
		数据模式、来源标签和评测输入清单 &
		流程测试 \\
		
		状态服务异常 &
		健康接口返回HTTP 503与\texttt{overall=degraded} &
		依赖状态、健康响应和容器日志 &
		部署测试 \\
		
		密钥或未授权访问风险 &
		密钥仅注入API容器，采用HTTPS、CORS白名单、JWT和配置隔离 &
		前端构建产物、容器配置和访问日志 &
		配置审查 \\
		
		高并发与持续限流 &
		沿用缓存、限速和任务预算，并监测尾部时延与错误恢复 &
		p50/p95时延、错误率、吞吐量和恢复时间 &
		待扩展验证 \\
		\bottomrule
	\end{tabularx}
\end{table}

\subsubsection{扩展验证计划}

现有测试已经覆盖来源失败、模型降级、缺失字段、精排异常和健康状态，但单元测试与小规模在线样本不足以刻画持续高负载和组合故障。后续可靠性验证应围绕三类场景展开。

第一，开展\textbf{并发与限流压力测试}。在固定查询集和冻结配置下逐级增加并发数，报告吞吐量、缓存命中率、p50/p95端到端时延、HTTP错误率和任务完成率，确定OpenAlex与Semantic Scholar限流条件下的稳定工作区间。

第二，开展\textbf{故障注入与恢复测试}。分别模拟单一学术API超时、模型服务不可达、Cross Encoder加载失败、Redis重启和Neo4j不可用，检查健康状态、降级路径、错误传播和恢复时间是否与设计一致；组合故障应单独记录，避免以单点测试代替系统级韧性结论。

第三，开展\textbf{跨学科与跨模型扩展验证}。CQC与EGRR需要在计算机科学、生物医学和交叉学科查询中比较约束解析、缺口触发和调用成本；模型规划接口应在相同查询、相同输出模式和相同Token上限下接入DeepSeek、Qwen、OpenAI兼容模型及其他可配置服务，比较结构化计划成功率、槽位F1、平均Token和降级比例。扩展验证的目标不是证明所有模型等价，而是确认系统接口、约束模式和降级语义能够跨模型复用。

\subsubsection{可靠性结论与能力边界}

ESASR已在数据源访问、模型规划、本地精排、元数据传播、状态服务和公网部署之间建立分层故障边界。外部依赖异常不会被隐藏为“零结果”，模型或精排不可用也不会直接扩大为完整链路失败；来源、缺失字段、调用成本和降级原因均能够进入统一运行轨迹。该设计使系统具备\textbf{可观察、可降级、可恢复和可审计}的工程基础。
\subsection{计划与分工}

项目以“责任主体、验收证据、版本节点”三者绑定的治理方式控制该风险：队长对需求基线、评测口径、风险闭环和最终验收负责；队员A对 CQC、EGRR、MEFR、LCER、CARS 及其量化证据负责；队员B对数据源、后端接口、前端联调、部署测试和提交材料负责。表中日期仅界定责任窗口，技术结论仍以冻结实验、自动测试和可运行产物为验收依据，从而保证工作分工服务于证据闭环而非形成过程流水账。

\begin{table}[H]
\centering
\small
\caption{项目阶段、版本节点与匿名成员分工}
\begin{tabularx}{\textwidth}{C{2.2cm} C{2.7cm} C{1.5cm} Y}
\toprule
\textbf{阶段} & \textbf{起止日期} & \textbf{版本节点} & \textbf{任务分工与验收产物} \\
\midrule
需求与架构 & 2026.05.22--2026.06.12 & v0.1--v0.2 & 难点：复合查询缺少统一验收口径。队长负责需求与指标边界，队员A负责约束模式，队员B负责架构与接口；以需求基线、数据流和接口规范作为验收证据。 \\
核心模块开发 & 2026.06.13--2026.07.10 & v0.3 & 难点：检索决策、数据源和排序模块需解耦。队员A负责 CQC、EGRR、MEFR、LCER 与 CARS，队员B负责双源接入、去重与缓存，队长负责接口一致性；以可运行检索链路验收。 \\
实验验证 & 2026.07.11--2026.08.08 & v0.4 & 难点：组件收益易受数据与输出规模混淆。队长负责冻结协议，队员A负责消融、置信区间与错误分析，队员B负责脚本和证据归档；以可复现实验包验收。 \\
集成与优化 & 2026.08.09--2026.08.26 & v0.5 & 难点：外部依赖故障可能破坏展示可信度。队员A负责性能与故障回归，队员B负责联调、部署和交互，队长负责风险闭环；以稳定演示系统和降级证据验收。 \\
提交定版 & 2026.08.27--2026.08.31 & v1.0 & 难点：代码、指标与材料必须保持一致。队长负责版本冻结，队员A负责数据与复现入口，队员B负责 PDF、附件和仓库；以一致性核验通过的提交包验收。 \\
\bottomrule
\end{tabularx}
\end{table}

\subsubsection{交付与复现}

项目复现要求环境、配置、数据划分和模型状态保持一致。交付包通过根目录启动脚本与Docker Compose统一构建服务，使用环境变量示例和配置模板隔离真实密钥，并通过独立开关控制本地Cross Encoder。模型未安装或精排异常时，系统保留MEFR基础排序并记录降级原因，避免将环境差异误判为算法变化。

实验目录保存Gold Set、逐查询预测、开发/测试拆分、指标与置信区间，同时记录数据哈希、随机种子、模型名称和阈值搜索结果。后端与CLI测试、前端构建、冻结实验、在线运行轨迹和服务健康状态共同构成验收依据，使汇总指标能够追溯至具体查询、候选和运行配置。

主报告只保留最终技术路线、必要消融和结论边界；中间方案、负结果、技术瓶颈及解决依据收录于《技术迭代与问题解决记录》；启动命令、API Key配置、目录结构和故障排查收录于《项目使用与复现说明》。这种分层结构既保持主报告的技术主线，也为关键决策和实验结果提供可查证的复现材料。

\clearpage
\section{参考资料}
{\small\begin{enumerate}
  \item 中国研究生创新实践系列大赛，第八届中国研究生人工智能创新大赛赛题揭榜公告，赛题“科研场景下复杂学术查询的智能论文搜索与推荐”，2026。
  \item Ajith A., Xia M., Chevalier A., et al. LitSearch: A Retrieval Benchmark for Scientific Literature Search. EMNLP, 2024, pp. 15068--15083.
  \item Jeong S., Baek J., Cho S., et al. Adaptive-RAG: Learning to Adapt Retrieval-Augmented Large Language Models through Question Complexity. NAACL, 2024.
  \item Asai A., Wu Z., Wang Y., et al. Self-RAG: Learning to Retrieve, Generate, and Critique through Self-Reflection. ICLR, 2024.
  \item Kang S., Zhang Y., Jiang P., et al. Taxonomy-guided Semantic Indexing for Academic Paper Search. EMNLP, 2024.
  \item Cormack G. V., Clarke C. L. A., Büttcher S. Reciprocal Rank Fusion Outperforms Condorcet and Individual Rank Learning Methods. SIGIR, 2009.
  \item Nogueira R., Cho K. Passage Re-ranking with BERT. arXiv:1901.04085, 2019.
  \item OpenAlex Documentation. \href{https://docs.openalex.org/}{https://docs.openalex.org/}.
  \item Semantic Scholar Academic Graph API Documentation. \href{https://api.semanticscholar.org/api-docs/}{https://api.semanticscholar.org/api-docs/}.
  \item FastAPI Documentation; Vue.js Documentation; Docker Compose Documentation.
\end{enumerate}}
\end{document}
